%============================================================
%  Bd. 02 - Theoreme der Logik: Vertiefte Beweise und Methoden
%============================================================

\documentclass[main.tex]{subfiles}

\title{CLASSICAL FIRST-ORDER LOGIC WITH EQUALITY}
\author{Martin Kunze}
\date{}

\ifSubfilesClassLoaded{
    \usepackage{xr-hyper}
  \externaldocument{registry/_B01}
}{
   % Code für als Subfile eingebunden
}

% Band-2 Registry/Debug unter festem Namen schreiben
\ifSubfilesClassLoaded{
  \directlua{
    thmlookup.registry_path = "registry/_B02.registry.tsv"
    thmlookup.debug_path    = "registry/_B02.debug.log"
    thmlookup.prepare_run()
    thmlookup.load_registry_file("registry/_B01.registry.tsv")
  }
}{}


\begin{document}
\maketitle
\tableofcontents
%\listoftheorems
\setcounter{file}{2}
\renewcommand{\FormulaBandID}{2}


\chapter{Pre-logic}

\section{Inferences for assisting proof development}

\begin{remark}
Die folgenden Hilfsinferenzregeln dienen nur der technischen Entwicklung von Beweisen in
Metamath. Sie gehören nicht zum eigentlichen mathematischen Kern des Kalküls und treten
in ausgearbeiteten Endbeweisen normalerweise nicht mehr auf.
\end{remark}

\FormulaThm[Identity inference for proof development]{%
  P \vdash P%
}{idi}{%
  \DeltaRow{Propositions}{P}%
}
\begin{tabproof}
  \proofstep{1}{P}{\rA}
\end{tabproof}

\FormulaThm[Auxiliary inference for proof development]{%
  P \dsep Q \vdash P%
}{a1ii}{%
  \DeltaRow{Propositions}{P\dsep Q}%
}
\begin{tabproof}
  \proofstep{1}{P}{\rA}
  \proofstep{2}{Q}{\rA}
  \proofstep{1,2}{P}{\rAEa{1}}
\end{tabproof}

\begin{hint}
In Metamath werden \FormulaRefAuto{idi} und \FormulaRefAuto{a1ii} als rein technische
Hilfsmittel zur Beweiskonstruktion verwendet. Sie können in einem mathematischen Skript
übergangen oder in einen Anhang ausgelagert werden.
\end{hint}

\chapter{Propositional calculus}

\section{Primitive well-formed formulas}

\FormulaDef[Negation as a formation rule]{%
  \text{If } P \text{ is a wff, then } \neg P \text{ is a wff.}%
}{wn}{%
  \DeltaRow{Propositions}{P}%
}

\FormulaDef[Implication as a formation rule]{%
  \text{If } P \text{ and } Q \text{ are wffs, then } (P \rightarrow Q) \text{ is a wff.}%
}{wi}{%
  \DeltaRow{Propositions}{P\dsep Q}%
}

\begin{remark}
Die Aussagenlogik behandelt allgemeine Wahrheiten über wohlgeformte Formeln unabhängig
davon, wie diese aufgebaut sind. Der einfachste Satz ist \((P \rightarrow P)\), der zwar
inhaltlich trivial erscheint, aber aus den Axiomen hergeleitet werden muss.
\end{remark}

\section{The axioms of propositional calculus}

\begin{remark}
Das hier verwendete Aussagenkalkül besteht aus den drei Grundaxiomen
\FormulaRefAuto{ax-1}, \FormulaRefAuto{ax-2}, \FormulaRefAuto{ax-3}
sowie der Schlussregel \FormulaRefAuto{ax-mp}. Die späteren Resultate des
Abschnitts über logische Implikation verwenden zunächst nur
\FormulaRefAuto{ax-mp}, \FormulaRefAuto{ax-1} und \FormulaRefAuto{ax-2}.
\end{remark}

\FormulaAxiom[Modus Ponens]{%
  P \dsep P \rightarrow Q \vdash Q%
}{ax-mp}{%
  \DeltaRow{Propositions}{P\dsep Q}%
}
\begin{tabproof}
  \proofstep{1}{P}{\rA}
  \proofstep{2}{P \rightarrow Q}{\rA}
  \proofstep{1,2}{Q}{\rRE{2,1}}
\end{tabproof}

\FormulaAxiom[Simp]{%
  Q \rightarrow (P \rightarrow Q)%
}{ax-1}{%
  \DeltaRow{Propositions}{P\dsep Q}%
}
\begin{tabproof}
  \proofstep{1}{Q}{\rA}
  \proofstep{2}{P}{\rA}
  \proofstep{1,2}{Q \land P}{\rAI{1,2}}
  \proofstep{1,2}{Q}{\rAEa{3}}
  \proofstep{1}{P \rightarrow Q}{\rRI{2,4}}
  \proofstep{}{Q \rightarrow (P \rightarrow Q)}{\rRI{1,5}}
\end{tabproof}

\FormulaAxiom[Frege]{%
  (P \rightarrow (Q \rightarrow R)) \rightarrow ((P \rightarrow Q) \rightarrow (P \rightarrow R))
}{ax-2}{%
  \DeltaRow{Propositions}{P\dsep Q\dsep R}
}
\begin{tabproof}
  \proofstep{1}{P \rightarrow (Q \rightarrow R)}{\rA}
  \proofstep{2}{P \rightarrow Q}{\rA}
  \proofstep{3}{P}{\rA}
  \proofstep{1,3}{Q \rightarrow R}{\rRE{1,3}}
  \proofstep{2,3}{Q}{\rRE{2,3}}
  \proofstep{1,2,3}{R}{\rRE{4,5}}
  \proofstep{1,2}{P \rightarrow R}{\rRI{3,6}}
  \proofstep{1}{(P \rightarrow Q) \rightarrow (P \rightarrow R)}{\rRI{2,7}}
  \proofstep{}{(P \rightarrow (Q \rightarrow R)) \rightarrow ((P \rightarrow Q) \rightarrow (P \rightarrow R))}{\rRI{1,8}}
\end{tabproof}

\FormulaThm[Transp]{%
  (\neg P \rightarrow \neg Q) \rightarrow (Q \rightarrow P)%
}{ax-3}{%
  \DeltaRow{Propositions}{P\dsep Q}%
}
\begin{tabproof}
  \proofstep{1}{\neg P \rightarrow \neg Q}{\rA}
  \proofstep{2}{Q}{\rA}
  \proofstep{3}{\neg P}{\rA}
  \proofstep{1,3}{\neg Q}{\rRE{1,3}}
  \proofstep{1,2,3}{\bot}{\rBI{2,4}}
  \proofstep{1,2}{P}{\rCE{3,5}}
  \proofstep{1}{Q \rightarrow P}{\rRI{2,7}}
  \proofstep{}{(\neg P \rightarrow \neg Q) \rightarrow (Q \rightarrow P)}{\rRI{1,8}}
\end{tabproof}

\section{Logical implication}

\FormulaThm[Double modus ponens]{%
  P \dsep Q \dsep P \rightarrow (Q \rightarrow R) \vdash R%
}{mp2}{%
  \DeltaRow{Propositions}{P\dsep Q\dsep R}%
}
\begin{tabproof}
  \proofstep{1}{P}{\rA}
  \proofstep{2}{Q}{\rA}
  \proofstep{3}{P \rightarrow (Q \rightarrow R)}{\rA}
  \proofstep{1,3}{Q \rightarrow R}{\FormulaRefAuto{ax-mp}{1,3}}
  \proofstep{1,2,3}{R}{\FormulaRefAuto{ax-mp}{2,4}}
\end{tabproof}

\FormulaThm[Double modus ponens]{%
  P \dsep P \rightarrow Q \dsep Q \rightarrow R \vdash R%
}{mp2b}{%
  \DeltaRow{Propositions}{P\dsep Q\dsep R}%
}
\begin{tabproof}
  \proofstep{1}{P}{\rA}
  \proofstep{2}{P \rightarrow Q}{\rA}
  \proofstep{3}{Q \rightarrow R}{\rA}
  \proofstep{1,2}{Q}{\FormulaRefAuto{ax-mp}{1,2}}
  \proofstep{1,2,3}{R}{\FormulaRefAuto{ax-mp}{4,3}}
\end{tabproof}

\FormulaThm[Inference introducing an antecedent]{%
  P \vdash Q \rightarrow P%
}{a1i}{%
  \DeltaRow{Propositions}{P\dsep Q}%
}
\begin{tabproof}
  \proofstep{1}{P}{\rA}
  \proofstep{}{P \rightarrow (Q \rightarrow P)}{\FormulaRefAuto{ax-1}}
  \proofstep{1}{Q \rightarrow P}{\FormulaRefAuto{ax-mp}{1,2}}
\end{tabproof}

\FormulaThm[Inference introducing two antecedents]{%
  P \vdash Q \rightarrow (R \rightarrow P)%
}{2a1i}{%
  \DeltaRow{Propositions}{P\dsep Q\dsep R}%
}
\begin{tabproof}
  \proofstep{1}{P}{\rA}
  \proofstep{1}{R \rightarrow P}{\FormulaRefAuto{a1i}{1}}
  \proofstep{1}{Q \rightarrow (R \rightarrow P)}{\FormulaRefAuto{a1i}{2}}
\end{tabproof}

\FormulaThm[Inference detaching an antecedent and introducing a new one]{%
  P \dsep P \rightarrow Q \vdash R \rightarrow Q%
}{mp1i}{%
  \DeltaRow{Propositions}{P\dsep Q\dsep R}%
}
\begin{tabproof}
  \proofstep{1}{P}{\rA}
  \proofstep{2}{P \rightarrow Q}{\rA}
  \proofstep{1,2}{Q}{\FormulaRefAuto{ax-mp}{1,2}}
  \proofstep{1,2}{R \rightarrow Q}{\FormulaRefAuto{a1i}{3}}
\end{tabproof}

\FormulaThm[Inference distributing an antecedent]{%
  P \rightarrow (Q \rightarrow R) \vdash (P \rightarrow Q) \rightarrow (P \rightarrow R)%
}{a2i}{%
  \DeltaRow{Propositions}{P\dsep Q\dsep R}%
}
\begin{tabproof}
  \proofstep{1}{P \rightarrow (Q \rightarrow R)}{\rA}
  \proofstep{}{(P \rightarrow (Q \rightarrow R)) \rightarrow ((P \rightarrow Q) \rightarrow (P \rightarrow R))}{\FormulaRefAuto{ax-2}}
  \proofstep{1}{(P \rightarrow Q) \rightarrow (P \rightarrow R)}{\FormulaRefAuto{ax-mp}{1,2}}
\end{tabproof}

\FormulaThm[A modus ponens deduction]{%
  P \rightarrow Q \dsep P \rightarrow (Q \rightarrow R) \vdash P \rightarrow R%
}{mpd}{%
  \DeltaRow{Propositions}{P\dsep Q\dsep R}%
}
\begin{tabproof}
  \proofstep{1}{P \rightarrow Q}{\rA}
  \proofstep{2}{P \rightarrow (Q \rightarrow R)}{\rA}
  \proofstep{2}{(P \rightarrow Q) \rightarrow (P \rightarrow R)}{\FormulaRefAuto{a2i}{2}}
  \proofstep{1,2}{P \rightarrow R}{\FormulaRefAuto{ax-mp}{1,3}}
\end{tabproof}

\FormulaThm[Inference adding common antecedents in an implication]{%
  P \rightarrow Q \vdash (R \rightarrow P) \rightarrow (R \rightarrow Q)%
}{imim2i}{%
  \DeltaRow{Propositions}{P\dsep Q\dsep R}%
}
\begin{tabproof}
  \proofstep{1}{P \rightarrow Q}{\rA}
  \proofstep{1}{R \rightarrow (P \rightarrow Q)}{\FormulaRefAuto{a1i}{1}}
  \proofstep{1}{(R \rightarrow P) \rightarrow (R \rightarrow Q)}{\FormulaRefAuto{a2i}{2}}
\end{tabproof}

\FormulaThm[Syllogism]{%
  P \rightarrow Q \dsep Q \rightarrow R \vdash P \rightarrow R%
}{syl}{%
  \DeltaRow{Propositions}{P\dsep Q\dsep R}%
}
\begin{tabproof}
  \proofstep{1}{P \rightarrow Q}{\rA}
  \proofstep{2}{Q \rightarrow R}{\rA}
  \proofstep{2}{P \rightarrow (Q \rightarrow R)}{\FormulaRefAuto{a1i}{2}}
  \proofstep{1,2}{P \rightarrow R}{\FormulaRefAuto{mpd}{1,3}}
\end{tabproof}

\FormulaThm[Inference chaining two syllogisms]{%
  P \rightarrow Q \dsep Q \rightarrow R \dsep R \rightarrow S \vdash P \rightarrow S%
}{3syl}{%
  \DeltaRow{Propositions}{P\dsep Q\dsep R\dsep S}%
}
\begin{tabproof}
  \proofstep{1}{P \rightarrow Q}{\rA}
  \proofstep{2}{Q \rightarrow R}{\rA}
  \proofstep{3}{R \rightarrow S}{\rA}
  \proofstep{1,2}{P \rightarrow R}{\FormulaRefAuto{syl}{1,2}}
  \proofstep{1,2,3}{P \rightarrow S}{\FormulaRefAuto{syl}{4,3}}
\end{tabproof}

\FormulaThm[Inference chaining three syllogisms]{%
  P \rightarrow Q \dsep Q \rightarrow R \dsep R \rightarrow S \dsep S \rightarrow T \vdash P \rightarrow T%
}{4syl}{%
  \DeltaRow{Propositions}{P\dsep Q\dsep R\dsep S\dsep T}%
}
\begin{tabproof}
  \proofstep{1}{P \rightarrow Q}{\rA}
  \proofstep{2}{Q \rightarrow R}{\rA}
  \proofstep{3}{R \rightarrow S}{\rA}
  \proofstep{4}{S \rightarrow T}{\rA}
  \proofstep{1,2,3}{P \rightarrow S}{\FormulaRefAuto{3syl}{1,2,3}}
  \proofstep{1,2,3,4}{P \rightarrow T}{\FormulaRefAuto{syl}{5,4}}
\end{tabproof}

\FormulaThm[A nested modus ponens inference]{%
  Q \dsep P \rightarrow (Q \rightarrow R) \vdash P \rightarrow R%
}{mpi}{%
  \DeltaRow{Propositions}{P\dsep Q\dsep R}%
}
\begin{tabproof}
  \proofstep{1}{Q}{\rA}
  \proofstep{2}{P \rightarrow (Q \rightarrow R)}{\rA}
  \proofstep{1}{P \rightarrow Q}{\FormulaRefAuto{a1i}{1}}
  \proofstep{1,2}{P \rightarrow R}{\FormulaRefAuto{mpd}{3,2}}
\end{tabproof}

\FormulaThm[A syllogism combined with a modus ponens inference]{%
  P \rightarrow Q \dsep R \dsep Q \rightarrow (R \rightarrow S) \vdash P \rightarrow S%
}{mpisyl}{%
  \DeltaRow{Propositions}{P\dsep Q\dsep R\dsep S}%
}
\begin{tabproof}
  \proofstep{1}{P \rightarrow Q}{\rA}
  \proofstep{2}{R}{\rA}
  \proofstep{3}{Q \rightarrow (R \rightarrow S)}{\rA}
  \proofstep{2,3}{Q \rightarrow S}{\FormulaRefAuto{mpi}{2,3}}
  \proofstep{1,2,3}{P \rightarrow S}{\FormulaRefAuto{syl}{1,4}}
\end{tabproof}

\FormulaThm[Principle of identity]{%
  P \rightarrow P%
}{id}{%
  \DeltaRow{Propositions}{P}%
}
\begin{tabproof}
  \proofstep{}{P \rightarrow (P \rightarrow P)}{\FormulaRefAuto{ax-1}}
  \proofstep{}{P \rightarrow ((P \rightarrow P) \rightarrow P)}{\FormulaRefAuto{ax-1}}
  \proofstep{}{P \rightarrow P}{\FormulaRefAuto{mpd}{1,2}}
\end{tabproof}
\begin{remark}
\FormulaRefAuto{id} ist die elementarste Tautologie des Aussagenkalküls. Obwohl die Aussage inhaltlich selbstverständlich wirkt, zeigt ihr Beweis, dass selbst die Selbstimplikation \(P \rightarrow P\) im syntaktischen Aufbau erst aus den Axiomen gewonnen werden muss. Die Variante \FormulaRefAuto{idALT} ist vor allem didaktisch interessant, weil sie denselben Satz unmittelbar aus den Grundaxiomen herleitet.
\end{remark}

\FormulaThm[Alternate proof of the principle of identity]{%
  P \rightarrow P'%
}{idALT}{%
  \DeltaRow{Propositions}{P}%
}
\begin{tabproof}
  \proofstep{}{P \rightarrow (P \rightarrow P)}{\FormulaRefAuto{ax-1}}
  \proofstep{}{P \rightarrow ((P \rightarrow P) \rightarrow P)}{\FormulaRefAuto{ax-1}}
  \proofstep{}{%
    (P \rightarrow ((P \rightarrow P) \rightarrow P)) \rightarrow
    ((P \rightarrow (P \rightarrow P)) \rightarrow (P \rightarrow P))
  }{\FormulaRefAuto{ax-2}}
  \proofstep{}{(P \rightarrow (P \rightarrow P)) \rightarrow (P \rightarrow P)}{\FormulaRefAuto{ax-mp}{2,3}}
  \proofstep{}{P \rightarrow P}{\FormulaRefAuto{ax-mp}{1,4}}
\end{tabproof}

\FormulaThm[Principle of identity with antecedent]{%
  P \rightarrow (Q \rightarrow Q)%
}{idd}{%
  \DeltaRow{Propositions}{P\dsep Q}%
}
\begin{tabproof}
  \proofstep{}{Q \rightarrow Q}{\FormulaRefAuto{id}}
  \proofstep{}{P \rightarrow (Q \rightarrow Q)}{\FormulaRefAuto{a1i}{1}}
\end{tabproof}

\FormulaThm[Deduction introducing an embedded antecedent]{%
  P \rightarrow Q \vdash P \rightarrow (R \rightarrow Q)%
}{a1d}{%
  \DeltaRow{Propositions}{P\dsep Q\dsep R}%
}
\begin{tabproof}
  \proofstep{1}{P \rightarrow Q}{\rA}
  \proofstep{}{Q \rightarrow (R \rightarrow Q)}{\FormulaRefAuto{ax-1}}
  \proofstep{1}{P \rightarrow (R \rightarrow Q)}{\FormulaRefAuto{syl}{1,2}}
\end{tabproof}

\FormulaThm[Deduction introducing two antecedents]{%
  P \rightarrow Q \vdash P \rightarrow (R \rightarrow (S \rightarrow Q))%
}{2a1d}{%
  \DeltaRow{Propositions}{P\dsep Q\dsep R\dsep S}%
}
\begin{tabproof}
  \proofstep{1}{P \rightarrow Q}{\rA}
  \proofstep{1}{P \rightarrow (S \rightarrow Q)}{\FormulaRefAuto{a1d}{1}}
  \proofstep{1}{P \rightarrow (R \rightarrow (S \rightarrow Q))}{\FormulaRefAuto{a1d}{2}}
\end{tabproof}

\FormulaThm[Add two antecedents to a wff]{%
  Q \rightarrow S \vdash P \rightarrow (Q \rightarrow (R \rightarrow S))%
}{a1i13}{%
  \DeltaRow{Propositions}{P\dsep Q\dsep R\dsep S}%
}
\begin{tabproof}
  \proofstep{1}{Q \rightarrow S}{\rA}
  \proofstep{1}{Q \rightarrow (R \rightarrow S)}{\FormulaRefAuto{a1d}{1}}
  \proofstep{1}{P \rightarrow (Q \rightarrow (R \rightarrow S))}{\FormulaRefAuto{a1i}{2}}
\end{tabproof}

\FormulaThm[A double form of ax-1]{%
  P \rightarrow (Q \rightarrow (R \rightarrow P))%
}{2a1}{%
  \DeltaRow{Propositions}{P\dsep Q\dsep R}%
}
\begin{tabproof}
  \proofstep{}{P \rightarrow P}{\FormulaRefAuto{id}}
  \proofstep{}{P \rightarrow (Q \rightarrow (R \rightarrow P))}{\FormulaRefAuto{2a1d}{1}}
\end{tabproof}

\FormulaThm[Deduction distributing an embedded antecedent]{%
  P \rightarrow (Q \rightarrow (R \rightarrow S)) \vdash
  P \rightarrow ((Q \rightarrow R) \rightarrow (Q \rightarrow S))%
}{a2d}{%
  \DeltaRow{Propositions}{P\dsep Q\dsep R\dsep S}%
}
\begin{tabproof}
  \proofstep{1}{P \rightarrow (Q \rightarrow (R \rightarrow S))}{\rA}
  \proofstep{}{(Q \rightarrow (R \rightarrow S)) \rightarrow ((Q \rightarrow R) \rightarrow (Q \rightarrow S))}{\FormulaRefAuto{ax-2}}
  \proofstep{1}{P \rightarrow ((Q \rightarrow R) \rightarrow (Q \rightarrow S))}{\FormulaRefAuto{syl}{1,2}}
\end{tabproof}

\FormulaThm[Syllogism inference with commutation of antecedents]{%
  P \rightarrow (Q \rightarrow R) \dsep Q \rightarrow (R \rightarrow S) \vdash P \rightarrow (Q \rightarrow S)%
}{sylcom}{%
  \DeltaRow{Propositions}{P\dsep Q\dsep R\dsep S}%
}
\begin{tabproof}
  \proofstep{1}{P \rightarrow (Q \rightarrow R)}{\rA}
  \proofstep{2}{Q \rightarrow (R \rightarrow S)}{\rA}
  \proofstep{2}{P \rightarrow (Q \rightarrow (R \rightarrow S))}{\FormulaRefAuto{a1i}{2}}
  \proofstep{2}{P \rightarrow ((Q \rightarrow R) \rightarrow (Q \rightarrow S))}{\FormulaRefAuto{a2d}{3}}
  \proofstep{1,2}{P \rightarrow (Q \rightarrow S)}{\FormulaRefAuto{mpd}{1,4}}
\end{tabproof}

\FormulaThm[Syllogism inference with commuted antecedents]{%
  P \rightarrow Q \dsep R \rightarrow (Q \rightarrow S) \vdash P \rightarrow (R \rightarrow S)%
}{syl5com}{%
  \DeltaRow{Propositions}{P\dsep Q\dsep R\dsep S}%
}
\begin{tabproof}
  \proofstep{1}{P \rightarrow Q}{\rA}
  \proofstep{2}{R \rightarrow (Q \rightarrow S)}{\rA}
  \proofstep{1}{P \rightarrow (R \rightarrow Q)}{\FormulaRefAuto{a1d}{1}}
  \proofstep{1,2}{P \rightarrow (R \rightarrow S)}{\FormulaRefAuto{sylcom}{3,2}}
\end{tabproof}

\FormulaThm[Inference that swaps antecedents in an implication]{%
  P \rightarrow (Q \rightarrow R) \vdash Q \rightarrow (P \rightarrow R)%
}{com12}{%
  \DeltaRow{Propositions}{P\dsep Q\dsep R}%
}
\begin{tabproof}
  \proofstep{}{Q \rightarrow Q}{\FormulaRefAuto{id}}
  \proofstep{1}{P \rightarrow (Q \rightarrow R)}{\rA}
  \proofstep{1}{Q \rightarrow (P \rightarrow R)}{\FormulaRefAuto{syl5com}{1,2}}
\end{tabproof}

\FormulaThm[A syllogism inference]{%
  P \rightarrow (Q \rightarrow R) \dsep S \rightarrow P \vdash Q \rightarrow (S \rightarrow R)%
}{syl11}{%
  \DeltaRow{Propositions}{P\dsep Q\dsep R\dsep S}%
}
\begin{tabproof}
  \proofstep{1}{P \rightarrow (Q \rightarrow R)}{\rA}
  \proofstep{2}{S \rightarrow P}{\rA}
  \proofstep{1}{Q \rightarrow (P \rightarrow R)}{\FormulaRefAuto{com12}{1}}
  \proofstep{1,2}{S \rightarrow (Q \rightarrow R)}{\FormulaRefAuto{syl5com}{2,3}}
  \proofstep{1,2}{Q \rightarrow (S \rightarrow R)}{\FormulaRefAuto{com12}{4}}
\end{tabproof}

\FormulaThm[A syllogism rule of inference]{%
  P \rightarrow Q \dsep R \rightarrow (Q \rightarrow S) \vdash R \rightarrow (P \rightarrow S)%
}{syl5}{%
  \DeltaRow{Propositions}{P\dsep Q\dsep R\dsep S}%
}
\begin{tabproof}
  \proofstep{1}{P \rightarrow Q}{\rA}
  \proofstep{2}{R \rightarrow (Q \rightarrow S)}{\rA}
  \proofstep{1,2}{P \rightarrow (R \rightarrow S)}{\FormulaRefAuto{syl5com}{1,2}}
  \proofstep{1,2}{R \rightarrow (P \rightarrow S)}{\FormulaRefAuto{com12}{3}}
\end{tabproof}

\FormulaThm[A syllogism rule of inference]{%
  P \rightarrow (Q \rightarrow R) \dsep R \rightarrow S \vdash P \rightarrow (Q \rightarrow S)%
}{syl6}{%
  \DeltaRow{Propositions}{P\dsep Q\dsep R\dsep S}%
}
\begin{tabproof}
  \proofstep{1}{P \rightarrow (Q \rightarrow R)}{\rA}
  \proofstep{2}{R \rightarrow S}{\rA}
  \proofstep{2}{Q \rightarrow (R \rightarrow S)}{\FormulaRefAuto{a1i}{2}}
  \proofstep{1,2}{P \rightarrow (Q \rightarrow S)}{\FormulaRefAuto{sylcom}{1,3}}
\end{tabproof}

\FormulaThm[Combine syl5 and syl6]{%
  P \rightarrow Q \dsep R \rightarrow (Q \rightarrow S) \dsep S \rightarrow T \vdash R \rightarrow (P \rightarrow T)%
}{syl56}{%
  \DeltaRow{Propositions}{P\dsep Q\dsep R\dsep S\dsep T}%
}
\begin{tabproof}
  \proofstep{1}{P \rightarrow Q}{\rA}
  \proofstep{2}{R \rightarrow (Q \rightarrow S)}{\rA}
  \proofstep{3}{S \rightarrow T}{\rA}
  \proofstep{1,2}{R \rightarrow (P \rightarrow S)}{\FormulaRefAuto{syl5}{1,2}}
  \proofstep{1,2,3}{R \rightarrow (P \rightarrow T)}{\FormulaRefAuto{syl6}{4,3}}
\end{tabproof}

\FormulaThm[Syllogism inference with commuted antecedents]{%
  P \rightarrow (Q \rightarrow R) \dsep R \rightarrow S \vdash Q \rightarrow (P \rightarrow S)%
}{syl6com}{%
  \DeltaRow{Propositions}{P\dsep Q\dsep R\dsep S}%
}
\begin{tabproof}
  \proofstep{1}{P \rightarrow (Q \rightarrow R)}{\rA}
  \proofstep{2}{R \rightarrow S}{\rA}
  \proofstep{1,2}{P \rightarrow (Q \rightarrow S)}{\FormulaRefAuto{syl6}{1,2}}
  \proofstep{1,2}{Q \rightarrow (P \rightarrow S)}{\FormulaRefAuto{com12}{3}}
\end{tabproof}

\FormulaThm[Modus ponens inference with commutation of antecedents]{%
  Q \rightarrow P \dsep P \rightarrow (Q \rightarrow R) \vdash Q \rightarrow R%
}{mpcom}{%
  \DeltaRow{Propositions}{P\dsep Q\dsep R}%
}
\begin{tabproof}
  \proofstep{1}{Q \rightarrow P}{\rA}
  \proofstep{2}{P \rightarrow (Q \rightarrow R)}{\rA}
  \proofstep{2}{Q \rightarrow (P \rightarrow R)}{\FormulaRefAuto{com12}{2}}
  \proofstep{1,2}{Q \rightarrow R}{\FormulaRefAuto{mpd}{1,3}}
\end{tabproof}

\FormulaThm[Syllogism inference with common nested antecedent]{%
  Q \rightarrow (P \rightarrow R) \dsep R \rightarrow (P \rightarrow S) \vdash Q \rightarrow (P \rightarrow S)%
}{syli}{%
  \DeltaRow{Propositions}{P\dsep Q\dsep R\dsep S}%
}
\begin{tabproof}
  \proofstep{1}{Q \rightarrow (P \rightarrow R)}{\rA}
  \proofstep{2}{R \rightarrow (P \rightarrow S)}{\rA}
  \proofstep{2}{P \rightarrow (R \rightarrow S)}{\FormulaRefAuto{com12}{2}}
  \proofstep{1,2}{Q \rightarrow (P \rightarrow S)}{\FormulaRefAuto{sylcom}{1,3}}
\end{tabproof}

\FormulaThm[Replace two antecedents]{%
  P \rightarrow Q \dsep R \rightarrow S \dsep Q \rightarrow (S \rightarrow T) \vdash P \rightarrow (R \rightarrow T)%
}{syl2im}{%
  \DeltaRow{Propositions}{P\dsep Q\dsep R\dsep S\dsep T}%
}
\begin{tabproof}
  \proofstep{1}{P \rightarrow Q}{\rA}
  \proofstep{2}{R \rightarrow S}{\rA}
  \proofstep{3}{Q \rightarrow (S \rightarrow T)}{\rA}
  \proofstep{2,3}{R \rightarrow (Q \rightarrow T)}{\FormulaRefAuto{syl5com}{2,3}}
  \proofstep{2,3}{Q \rightarrow (R \rightarrow T)}{\FormulaRefAuto{com12}{4}}
  \proofstep{1,2,3}{P \rightarrow (R \rightarrow T)}{\FormulaRefAuto{syl}{1,5}}
\end{tabproof}

\FormulaThm[A commuted version of syl2im]{%
  P \rightarrow Q \dsep R \rightarrow S \dsep Q \rightarrow (S \rightarrow T) \vdash R \rightarrow (P \rightarrow T)%
}{syl2imc}{%
  \DeltaRow{Propositions}{P\dsep Q\dsep R\dsep S\dsep T}%
}
\begin{tabproof}
  \proofstep{1}{P \rightarrow Q}{\rA}
  \proofstep{2}{R \rightarrow S}{\rA}
  \proofstep{3}{Q \rightarrow (S \rightarrow T)}{\rA}
  \proofstep{1,2,3}{P \rightarrow (R \rightarrow T)}{\FormulaRefAuto{syl2im}{1,2,3}}
  \proofstep{1,2,3}{R \rightarrow (P \rightarrow T)}{\FormulaRefAuto{com12}{4}}
\end{tabproof}

\FormulaThm[Assertion]{%
  P \rightarrow ((P \rightarrow Q) \rightarrow Q)%
}{pm2.27}{%
  \DeltaRow{Propositions}{P\dsep Q}%
}
\begin{tabproof}
  \proofstep{}{(P \rightarrow Q) \rightarrow (P \rightarrow Q)}{\FormulaRefAuto{id}}
  \proofstep{}{%
    ((P \rightarrow Q) \rightarrow (P \rightarrow Q)) \rightarrow
    (((P \rightarrow Q) \rightarrow P) \rightarrow ((P \rightarrow Q) \rightarrow Q))
  }{\FormulaRefAuto{ax-2}}
  \proofstep{}{((P \rightarrow Q) \rightarrow P) \rightarrow ((P \rightarrow Q) \rightarrow Q)}{\FormulaRefAuto{ax-mp}{1,2}}
  \proofstep{}{P \rightarrow ((P \rightarrow Q) \rightarrow P)}{\FormulaRefAuto{ax-1}}
  \proofstep{}{%
    (P \rightarrow ((P \rightarrow Q) \rightarrow P)) \rightarrow
    (P \rightarrow ((P \rightarrow Q) \rightarrow Q))
  }{\FormulaRefAuto{imim2i}{3}}
  \proofstep{}{P \rightarrow ((P \rightarrow Q) \rightarrow Q)}{\FormulaRefAuto{ax-mp}{4,5}}
\end{tabproof}

\FormulaThm[A nested modus ponens deduction]{%
  P \rightarrow (Q \rightarrow R) \dsep P \rightarrow (Q \rightarrow (R \rightarrow S)) \vdash P \rightarrow (Q \rightarrow S)%
}{mpdd}{%
  \DeltaRow{Propositions}{P\dsep Q\dsep R\dsep S}%
}
\begin{tabproof}
  \proofstep{1}{P \rightarrow (Q \rightarrow R)}{\rA}
  \proofstep{2}{P \rightarrow (Q \rightarrow (R \rightarrow S))}{\rA}
  \proofstep{2}{P \rightarrow ((Q \rightarrow R) \rightarrow (Q \rightarrow S))}{\FormulaRefAuto{a2d}{2}}
  \proofstep{1,2}{P \rightarrow (Q \rightarrow S)}{\FormulaRefAuto{mpd}{1,3}}
\end{tabproof}

\FormulaThm[A nested modus ponens deduction]{%
  P \rightarrow R \dsep P \rightarrow (Q \rightarrow (R \rightarrow S)) \vdash P \rightarrow (Q \rightarrow S)%
}{mpid}{%
  \DeltaRow{Propositions}{P\dsep Q\dsep R\dsep S}%
}
\begin{tabproof}
  \proofstep{1}{P \rightarrow R}{\rA}
  \proofstep{2}{P \rightarrow (Q \rightarrow (R \rightarrow S))}{\rA}
  \proofstep{1}{P \rightarrow (Q \rightarrow R)}{\FormulaRefAuto{a1d}{1}}
  \proofstep{1,2}{P \rightarrow (Q \rightarrow S)}{\FormulaRefAuto{mpdd}{3,2}}
\end{tabproof}

\FormulaThm[A nested modus ponens deduction]{%
  Q \rightarrow R \dsep P \rightarrow (Q \rightarrow (R \rightarrow S)) \vdash P \rightarrow (Q \rightarrow S)%
}{mpdi}{%
  \DeltaRow{Propositions}{P\dsep Q\dsep R\dsep S}%
}
\begin{tabproof}
  \proofstep{1}{Q \rightarrow R}{\rA}
  \proofstep{2}{P \rightarrow (Q \rightarrow (R \rightarrow S))}{\rA}
  \proofstep{1}{P \rightarrow (Q \rightarrow R)}{\FormulaRefAuto{a1i}{1}}
  \proofstep{1,2}{P \rightarrow (Q \rightarrow S)}{\FormulaRefAuto{mpdd}{3,2}}
\end{tabproof}

\FormulaThm[A doubly nested modus ponens inference]{%
  R \dsep P \rightarrow (Q \rightarrow (R \rightarrow S)) \vdash P \rightarrow (Q \rightarrow S)%
}{mpii}{%
  \DeltaRow{Propositions}{P\dsep Q\dsep R\dsep S}%
}
\begin{tabproof}
  \proofstep{1}{R}{\rA}
  \proofstep{2}{P \rightarrow (Q \rightarrow (R \rightarrow S))}{\rA}
  \proofstep{1}{Q \rightarrow R}{\FormulaRefAuto{a1i}{1}}
  \proofstep{1,2}{P \rightarrow (Q \rightarrow S)}{\FormulaRefAuto{mpdi}{3,2}}
\end{tabproof}

\FormulaThm[Syllogism deduction]{%
  P \rightarrow (Q \rightarrow R) \dsep P \rightarrow (R \rightarrow S) \vdash P \rightarrow (Q \rightarrow S)%
}{syld}{%
  \DeltaRow{Propositions}{P\dsep Q\dsep R\dsep S}%
}
\begin{tabproof}
  \proofstep{1}{P \rightarrow (Q \rightarrow R)}{\rA}
  \proofstep{2}{P \rightarrow (R \rightarrow S)}{\rA}
  \proofstep{2}{P \rightarrow (Q \rightarrow (R \rightarrow S))}{\FormulaRefAuto{a1d}{2}}
  \proofstep{1,2}{P \rightarrow (Q \rightarrow S)}{\FormulaRefAuto{mpdd}{1,3}}
\end{tabproof}

\FormulaThm[Syllogism deduction]{%
  P \rightarrow (Q \rightarrow R) \dsep P \rightarrow (R \rightarrow S) \vdash Q \rightarrow (P \rightarrow S)%
}{syldc}{%
  \DeltaRow{Propositions}{P\dsep Q\dsep R\dsep S}%
}
\begin{tabproof}
  \proofstep{1}{P \rightarrow (Q \rightarrow R)}{\rA}
  \proofstep{2}{P \rightarrow (R \rightarrow S)}{\rA}
  \proofstep{1,2}{P \rightarrow (Q \rightarrow S)}{\FormulaRefAuto{syld}{1,2}}
  \proofstep{1,2}{Q \rightarrow (P \rightarrow S)}{\FormulaRefAuto{com12}{3}}
\end{tabproof}

\FormulaThm[A double modus ponens deduction]{%
  P \rightarrow Q \dsep P \rightarrow R \dsep P \rightarrow (Q \rightarrow (R \rightarrow S)) \vdash P \rightarrow S%
}{mp2d}{%
  \DeltaRow{Propositions}{P\dsep Q\dsep R\dsep S}%
}
\begin{tabproof}
  \proofstep{1}{P \rightarrow Q}{\rA}
  \proofstep{2}{P \rightarrow R}{\rA}
  \proofstep{3}{P \rightarrow (Q \rightarrow (R \rightarrow S))}{\rA}
  \proofstep{2,3}{P \rightarrow (Q \rightarrow S)}{\FormulaRefAuto{mpid}{2,3}}
  \proofstep{1,2,3}{P \rightarrow S}{\FormulaRefAuto{mpd}{1,4}}
\end{tabproof}

\FormulaThm[Double deduction introducing an antecedent]{%
  P \rightarrow (Q \rightarrow R) \vdash P \rightarrow (Q \rightarrow (S \rightarrow R))%
}{a1dd}{%
  \DeltaRow{Propositions}{P\dsep Q\dsep R\dsep S}%
}
\begin{tabproof}
  \proofstep{1}{P \rightarrow (Q \rightarrow R)}{\rA}
  \proofstep{}{R \rightarrow (S \rightarrow R)}{\FormulaRefAuto{ax-1}}
  \proofstep{1}{P \rightarrow (Q \rightarrow (S \rightarrow R))}{\FormulaRefAuto{syl6}{1,2}}
\end{tabproof}

\FormulaThm[Double deduction introducing two antecedents]{%
  P \rightarrow (Q \rightarrow R) \vdash P \rightarrow (Q \rightarrow (S \rightarrow (T \rightarrow R)))%
}{2a1dd}{%
  \DeltaRow{Propositions}{P\dsep Q\dsep R\dsep S\dsep T}%
}
\begin{tabproof}
  \proofstep{1}{P \rightarrow (Q \rightarrow R)}{\rA}
  \proofstep{1}{P \rightarrow (Q \rightarrow (T \rightarrow R))}{\FormulaRefAuto{a1dd}{1}}
  \proofstep{1}{P \rightarrow (Q \rightarrow (S \rightarrow (T \rightarrow R)))}{\FormulaRefAuto{a1dd}{2}}
\end{tabproof}

\FormulaThm[Inference absorbing redundant antecedent]{%
  P \rightarrow (P \rightarrow Q) \vdash P \rightarrow Q%
}{pm2.43i}{%
  \DeltaRow{Propositions}{P\dsep Q}%
}
\begin{tabproof}
  \proofstep{1}{P \rightarrow (P \rightarrow Q)}{\rA}
  \proofstep{}{(P \rightarrow (P \rightarrow Q)) \rightarrow (P \rightarrow Q)}{\FormulaRefAuto{pm2.43}}
  \proofstep{1}{P \rightarrow Q}{\FormulaRefAuto{ax-mp}{1,2}}
\end{tabproof}

\FormulaThm[Deduction absorbing redundant antecedent]{%
  P \rightarrow (Q \rightarrow (Q \rightarrow R)) \vdash P \rightarrow (Q \rightarrow R)%
}{pm2.43d}{%
  \DeltaRow{Propositions}{P\dsep Q\dsep R}%
}
\begin{tabproof}
  \proofstep{}{Q \rightarrow Q}{\FormulaRefAuto{id}}
  \proofstep{1}{P \rightarrow (Q \rightarrow (Q \rightarrow R))}{\rA}
  \proofstep{1}{P \rightarrow (Q \rightarrow R)}{\FormulaRefAuto{mpdi}{1,2}}
\end{tabproof}

\FormulaThm[Inference absorbing redundant antecedent]{%
  Q \rightarrow (P \rightarrow (Q \rightarrow R)) \vdash Q \rightarrow (P \rightarrow R)%
}{pm2.43a}{%
  \DeltaRow{Propositions}{P\dsep Q\dsep R}%
}
\begin{tabproof}
  \proofstep{}{Q \rightarrow Q}{\FormulaRefAuto{id}}
  \proofstep{1}{Q \rightarrow (P \rightarrow (Q \rightarrow R))}{\rA}
  \proofstep{1}{Q \rightarrow (P \rightarrow R)}{\FormulaRefAuto{mpid}{2,1}}
\end{tabproof}

\FormulaThm[Inference absorbing redundant antecedent]{%
  Q \rightarrow (P \rightarrow (Q \rightarrow R)) \vdash P \rightarrow (Q \rightarrow R)%
}{pm2.43b}{%
  \DeltaRow{Propositions}{P\dsep Q\dsep R}%
}
\begin{tabproof}
  \proofstep{1}{Q \rightarrow (P \rightarrow (Q \rightarrow R))}{\rA}
  \proofstep{1}{Q \rightarrow (P \rightarrow R)}{\FormulaRefAuto{pm2.43a}{1}}
  \proofstep{1}{P \rightarrow (Q \rightarrow R)}{\FormulaRefAuto{com12}{2}}
\end{tabproof}

\FormulaThm[Absorption of redundant antecedent]{%
  (P \rightarrow (P \rightarrow Q)) \rightarrow (P \rightarrow Q)%
}{pm2.43}{%
  \DeltaRow{Propositions}{P\dsep Q}%
}
\begin{tabproof}
  \proofstep{}{P \rightarrow ((P \rightarrow Q) \rightarrow Q)}{\FormulaRefAuto{pm2.27}}
  \proofstep{}{(P \rightarrow ((P \rightarrow Q) \rightarrow Q)) \rightarrow ((P \rightarrow (P \rightarrow Q)) \rightarrow (P \rightarrow Q))}{\FormulaRefAuto{a2i}{1}}
  \proofstep{}{(P \rightarrow (P \rightarrow Q)) \rightarrow (P \rightarrow Q)}{\FormulaRefAuto{ax-mp}{1,2}}
\end{tabproof}

\FormulaThm[Deduction adding nested antecedents]{%
  P \rightarrow (Q \rightarrow R) \vdash P \rightarrow ((S \rightarrow Q) \rightarrow (S \rightarrow R))%
}{imim2d}{%
  \DeltaRow{Propositions}{P\dsep Q\dsep R\dsep S}%
}
\begin{tabproof}
  \proofstep{1}{P \rightarrow (Q \rightarrow R)}{\rA}
  \proofstep{1}{P \rightarrow (S \rightarrow (Q \rightarrow R))}{\FormulaRefAuto{a1d}{1}}
  \proofstep{1}{P \rightarrow ((S \rightarrow Q) \rightarrow (S \rightarrow R))}{\FormulaRefAuto{a2d}{2}}
\end{tabproof}

\FormulaThm[A closed form of syllogism]{%
  (P \rightarrow Q) \rightarrow ((R \rightarrow P) \rightarrow (R \rightarrow Q))%
}{imim2}{%
  \DeltaRow{Propositions}{P\dsep Q\dsep R}%
}
\begin{tabproof}
  \proofstep{}{(P \rightarrow Q) \rightarrow (R \rightarrow (P \rightarrow Q))}{\FormulaRefAuto{ax-1}}
  \proofstep{}{(P \rightarrow Q) \rightarrow ((R \rightarrow P) \rightarrow (R \rightarrow Q))}{\FormulaRefAuto{a2d}{1}}
\end{tabproof}

\FormulaThm[Deduction embedding an antecedent]{%
  P \rightarrow Q \dsep P \rightarrow (R \rightarrow S) \vdash P \rightarrow ((Q \rightarrow R) \rightarrow S)%
}{embantd}{%
  \DeltaRow{Propositions}{P\dsep Q\dsep R\dsep S}%
}
\begin{tabproof}
  \proofstep{1}{P \rightarrow Q}{\rA}
  \proofstep{2}{P \rightarrow (R \rightarrow S)}{\rA}
  \proofstep{2}{P \rightarrow ((Q \rightarrow R) \rightarrow (Q \rightarrow S))}{\FormulaRefAuto{imim2d}{2}}
  \proofstep{1,2}{P \rightarrow ((Q \rightarrow R) \rightarrow S)}{\FormulaRefAuto{mpid}{1,3}}
\end{tabproof}

\FormulaThm[Triple syllogism deduction]{%
  P \rightarrow (Q \rightarrow R) \dsep P \rightarrow (R \rightarrow S) \dsep P \rightarrow (S \rightarrow T) \vdash P \rightarrow (Q \rightarrow T)%
}{3syld}{%
  \DeltaRow{Propositions}{P\dsep Q\dsep R\dsep S\dsep T}%
}
\begin{tabproof}
  \proofstep{1}{P \rightarrow (Q \rightarrow R)}{\rA}
  \proofstep{2}{P \rightarrow (R \rightarrow S)}{\rA}
  \proofstep{1,2}{P \rightarrow (Q \rightarrow S)}{\FormulaRefAuto{syld}{1,2}}
  \proofstep{4}{P \rightarrow (S \rightarrow T)}{\rA}
  \proofstep{1,2,4}{P \rightarrow (Q \rightarrow T)}{\FormulaRefAuto{syld}{3,4}}
\end{tabproof}

\FormulaThm[A double syllogism inference]{%
  P \rightarrow Q \dsep P \rightarrow (R \rightarrow S) \dsep Q \rightarrow (S \rightarrow T) \vdash P \rightarrow (R \rightarrow T)%
}{sylsyld}{%
  \DeltaRow{Propositions}{P\dsep Q\dsep R\dsep S\dsep T}%
}
\begin{tabproof}
  \proofstep{1}{P \rightarrow Q}{\rA}
  \proofstep{2}{P \rightarrow (R \rightarrow S)}{\rA}
  \proofstep{3}{Q \rightarrow (S \rightarrow T)}{\rA}
  \proofstep{1,3}{P \rightarrow (S \rightarrow T)}{\FormulaRefAuto{syl}{1,3}}
  \proofstep{1,2,3}{P \rightarrow (R \rightarrow T)}{\FormulaRefAuto{syld}{2,4}}
\end{tabproof}

\FormulaThm[Inference joining two implications]{%
  P \rightarrow Q \dsep R \rightarrow S \vdash (Q \rightarrow R) \rightarrow (P \rightarrow S)%
}{imim12i}{%
  \DeltaRow{Propositions}{P\dsep Q\dsep R\dsep S}%
}
\begin{tabproof}
  \proofstep{1}{P \rightarrow Q}{\rA}
  \proofstep{2}{R \rightarrow S}{\rA}
  \proofstep{2}{(Q \rightarrow R) \rightarrow (Q \rightarrow S)}{\FormulaRefAuto{imim2i}{2}}
  \proofstep{1,2}{(Q \rightarrow R) \rightarrow (P \rightarrow S)}{\FormulaRefAuto{syl5}{1,3}}
\end{tabproof}

\FormulaThm[Inference adding common consequents in an implication]{%
  P \rightarrow Q \vdash (Q \rightarrow R) \rightarrow (P \rightarrow R)%
}{imim1i}{%
  \DeltaRow{Propositions}{P\dsep Q\dsep R}%
}
\begin{tabproof}
  \proofstep{1}{P \rightarrow Q}{\rA}
  \proofstep{}{R \rightarrow R}{\FormulaRefAuto{id}}
  \proofstep{1}{(Q \rightarrow R) \rightarrow (P \rightarrow R)}{\FormulaRefAuto{imim12i}{1,2}}
\end{tabproof}

\FormulaThm[Inference adding three nested antecedents]{%
  P \rightarrow (Q \rightarrow R) \vdash (S \rightarrow P) \rightarrow ((S \rightarrow Q) \rightarrow (S \rightarrow R))%
}{imim3i}{%
  \DeltaRow{Propositions}{P\dsep Q\dsep R\dsep S}%
}
\begin{tabproof}
  \proofstep{1}{P \rightarrow (Q \rightarrow R)}{\rA}
  \proofstep{1}{(S \rightarrow P) \rightarrow (S \rightarrow (Q \rightarrow R))}{\FormulaRefAuto{imim2i}{1}}
  \proofstep{1}{(S \rightarrow P) \rightarrow ((S \rightarrow Q) \rightarrow (S \rightarrow R))}{\FormulaRefAuto{a2d}{2}}
\end{tabproof}


\FormulaThm[A syllogism inference combined with contraction]{%
  P \rightarrow Q \dsep P \rightarrow R \dsep Q \rightarrow (R \rightarrow S) \vdash P \rightarrow S%
}{sylc}{%
  \DeltaRow{Propositions}{P\dsep Q\dsep R\dsep S}%
}
\begin{tabproof}
  \proofstep{1}{P \rightarrow Q}{\rA}
  \proofstep{2}{P \rightarrow R}{\rA}
  \proofstep{3}{Q \rightarrow (R \rightarrow S)}{\rA}
  \proofstep{1,2,3}{P \rightarrow (P \rightarrow S)}{\FormulaRefAuto{syl2im}{1,2,3}}
  \proofstep{1,2,3}{P \rightarrow S}{\FormulaRefAuto{pm2.43i}{4}}
\end{tabproof}

\FormulaThm[A syllogism inference combined with contraction]{%
  P \rightarrow Q \dsep P \rightarrow R \dsep P \rightarrow S \dsep Q \rightarrow (R \rightarrow (S \rightarrow T)) \vdash P \rightarrow T%
}{syl3c}{%
  \DeltaRow{Propositions}{P\dsep Q\dsep R\dsep S\dsep T}%
}
\begin{tabproof}
  \proofstep{1}{P \rightarrow Q}{\rA}
  \proofstep{2}{P \rightarrow R}{\rA}
  \proofstep{3}{P \rightarrow S}{\rA}
  \proofstep{4}{Q \rightarrow (R \rightarrow (S \rightarrow T))}{\rA}
  \proofstep{1,2,4}{P \rightarrow (S \rightarrow T)}{\FormulaRefAuto{sylc}{1,2,4}}
  \proofstep{1,2,3,4}{P \rightarrow T}{\FormulaRefAuto{mpd}{3,5}}
\end{tabproof}

\FormulaThm[A syllogism inference]{%
  P \rightarrow (Q \rightarrow R) \dsep S \dsep R \rightarrow (S \rightarrow T) \vdash P \rightarrow (Q \rightarrow T)%
}{syl6mpi}{%
  \DeltaRow{Propositions}{P\dsep Q\dsep R\dsep S\dsep T}%
}
\begin{tabproof}
  \proofstep{1}{P \rightarrow (Q \rightarrow R)}{\rA}
  \proofstep{2}{S}{\rA}
  \proofstep{3}{R \rightarrow (S \rightarrow T)}{\rA}
  \proofstep{2,3}{R \rightarrow T}{\FormulaRefAuto{mpi}{2,3}}
  \proofstep{1,2,3}{P \rightarrow (Q \rightarrow T)}{\FormulaRefAuto{syl6}{1,4}}
\end{tabproof}

\FormulaThm[Modus ponens combined with a syllogism inference]{%
  P \dsep Q \rightarrow R \dsep P \rightarrow (R \rightarrow S) \vdash Q \rightarrow S%
}{mpsyl}{%
  \DeltaRow{Propositions}{P\dsep Q\dsep R\dsep S}%
}
\begin{tabproof}
  \proofstep{1}{P}{\rA}
  \proofstep{2}{Q \rightarrow R}{\rA}
  \proofstep{3}{P \rightarrow (R \rightarrow S)}{\rA}
  \proofstep{1}{Q \rightarrow P}{\FormulaRefAuto{a1i}{1}}
  \proofstep{1,2,3}{Q \rightarrow S}{\FormulaRefAuto{sylc}{4,2,3}}
\end{tabproof}

\FormulaThm[Modus ponens combined with a double syllogism inference]{%
  P \dsep Q \rightarrow (R \rightarrow S) \dsep P \rightarrow (S \rightarrow T) \vdash Q \rightarrow (R \rightarrow T)%
}{mpsylsyld}{%
  \DeltaRow{Propositions}{P\dsep Q\dsep R\dsep S\dsep T}%
}
\begin{tabproof}
  \proofstep{1}{P}{\rA}
  \proofstep{2}{Q \rightarrow (R \rightarrow S)}{\rA}
  \proofstep{3}{P \rightarrow (S \rightarrow T)}{\rA}
  \proofstep{1}{Q \rightarrow P}{\FormulaRefAuto{a1i}{1}}
  \proofstep{1,2,3}{Q \rightarrow (R \rightarrow T)}{\FormulaRefAuto{sylsyld}{4,2,3}}
\end{tabproof}

\FormulaThm[Inference combining syl6 with contraction]{%
  P \rightarrow (Q \rightarrow R) \dsep P \rightarrow (Q \rightarrow S) \dsep R \rightarrow (S \rightarrow T) \vdash P \rightarrow (Q \rightarrow T)%
}{syl6c}{%
  \DeltaRow{Propositions}{P\dsep Q\dsep R\dsep S\dsep T}%
}
\begin{tabproof}
  \proofstep{1}{P \rightarrow (Q \rightarrow R)}{\rA}
  \proofstep{2}{P \rightarrow (Q \rightarrow S)}{\rA}
  \proofstep{3}{R \rightarrow (S \rightarrow T)}{\rA}
  \proofstep{1,3}{P \rightarrow (Q \rightarrow (S \rightarrow T))}{\FormulaRefAuto{syl6}{1,3}}
  \proofstep{1,2,3}{P \rightarrow (Q \rightarrow T)}{\FormulaRefAuto{mpdd}{2,4}}
\end{tabproof}

\FormulaThm[A syllogism inference combined with contraction]{%
  P \rightarrow (Q \rightarrow R) \dsep P \rightarrow S \dsep R \rightarrow (S \rightarrow T) \vdash P \rightarrow (Q \rightarrow T)%
}{syl6ci}{%
  \DeltaRow{Propositions}{P\dsep Q\dsep R\dsep S\dsep T}%
}
\begin{tabproof}
  \proofstep{1}{P \rightarrow (Q \rightarrow R)}{\rA}
  \proofstep{2}{P \rightarrow S}{\rA}
  \proofstep{2}{P \rightarrow (Q \rightarrow S)}{\FormulaRefAuto{a1d}{2}}
  \proofstep{4}{R \rightarrow (S \rightarrow T)}{\rA}
  \proofstep{1,2,4}{P \rightarrow (Q \rightarrow T)}{\FormulaRefAuto{syl6c}{1,3,4}}
\end{tabproof}

\FormulaThm[Nested syllogism deduction]{%
  P \rightarrow (Q \rightarrow (R \rightarrow S)) \dsep P \rightarrow (Q \rightarrow (S \rightarrow T)) \vdash P \rightarrow (Q \rightarrow (R \rightarrow T))%
}{syldd}{%
  \DeltaRow{Propositions}{P\dsep Q\dsep R\dsep S\dsep T}%
}
\begin{tabproof}
  \proofstep{1}{P \rightarrow (Q \rightarrow (R \rightarrow S))}{\rA}
  \proofstep{2}{P \rightarrow (Q \rightarrow (S \rightarrow T))}{\rA}
  \proofstep{}{(S \rightarrow T) \rightarrow ((R \rightarrow S) \rightarrow (R \rightarrow T))}{\FormulaRefAuto{imim2}}
  \proofstep{1,2}{P \rightarrow (Q \rightarrow (R \rightarrow T))}{\FormulaRefAuto{syl6c}{2,1,3}}
\end{tabproof}

\FormulaThm[A nested syllogism deduction]{%
  P \rightarrow (Q \rightarrow R) \dsep P \rightarrow (S \rightarrow (R \rightarrow T)) \vdash P \rightarrow (S \rightarrow (Q \rightarrow T))%
}{syl5d}{%
  \DeltaRow{Propositions}{P\dsep Q\dsep R\dsep S\dsep T}%
}
\begin{tabproof}
  \proofstep{1}{P \rightarrow (Q \rightarrow R)}{\rA}
  \proofstep{1}{P \rightarrow (S \rightarrow (Q \rightarrow R))}{\FormulaRefAuto{a1d}{1}}
  \proofstep{3}{P \rightarrow (S \rightarrow (R \rightarrow T))}{\rA}
  \proofstep{1,3}{P \rightarrow (S \rightarrow (Q \rightarrow T))}{\FormulaRefAuto{syldd}{2,3}}
\end{tabproof}

\FormulaThm[A syllogism rule of inference]{%
  P \rightarrow Q \dsep R \rightarrow (S \rightarrow (Q \rightarrow T)) \vdash R \rightarrow (S \rightarrow (P \rightarrow T))%
}{syl7}{%
  \DeltaRow{Propositions}{P\dsep Q\dsep R\dsep S\dsep T}%
}
\begin{tabproof}
  \proofstep{1}{P \rightarrow Q}{\rA}
  \proofstep{1}{R \rightarrow (P \rightarrow Q)}{\FormulaRefAuto{a1i}{1}}
  \proofstep{3}{R \rightarrow (S \rightarrow (Q \rightarrow T))}{\rA}
  \proofstep{1,3}{R \rightarrow (S \rightarrow (P \rightarrow T))}{\FormulaRefAuto{syl5d}{2,3}}
\end{tabproof}

\FormulaThm[A nested syllogism deduction]{%
  P \rightarrow (Q \rightarrow (R \rightarrow S)) \dsep P \rightarrow (S \rightarrow T) \vdash P \rightarrow (Q \rightarrow (R \rightarrow T))%
}{syl6d}{%
  \DeltaRow{Propositions}{P\dsep Q\dsep R\dsep S\dsep T}%
}
\begin{tabproof}
  \proofstep{1}{P \rightarrow (Q \rightarrow (R \rightarrow S))}{\rA}
  \proofstep{2}{P \rightarrow (S \rightarrow T)}{\rA}
  \proofstep{2}{P \rightarrow (Q \rightarrow (S \rightarrow T))}{\FormulaRefAuto{a1d}{2}}
  \proofstep{1,2}{P \rightarrow (Q \rightarrow (R \rightarrow T))}{\FormulaRefAuto{syldd}{1,3}}
\end{tabproof}

\FormulaThm[A syllogism rule of inference]{%
  P \rightarrow (Q \rightarrow (R \rightarrow S)) \dsep S \rightarrow T \vdash P \rightarrow (Q \rightarrow (R \rightarrow T))%
}{syl8}{%
  \DeltaRow{Propositions}{P\dsep Q\dsep R\dsep S\dsep T}%
}
\begin{tabproof}
  \proofstep{1}{P \rightarrow (Q \rightarrow (R \rightarrow S))}{\rA}
  \proofstep{2}{S \rightarrow T}{\rA}
  \proofstep{2}{P \rightarrow (S \rightarrow T)}{\FormulaRefAuto{a1i}{2}}
  \proofstep{1,2}{P \rightarrow (Q \rightarrow (R \rightarrow T))}{\FormulaRefAuto{syl6d}{1,3}}
\end{tabproof}

\FormulaThm[A nested syllogism inference with different antecedents]{%
  P \rightarrow (Q \rightarrow R) \dsep S \rightarrow (R \rightarrow T) \vdash P \rightarrow (S \rightarrow (Q \rightarrow T))%
}{syl9}{%
  \DeltaRow{Propositions}{P\dsep Q\dsep R\dsep S\dsep T}%
}
\begin{tabproof}
  \proofstep{1}{P \rightarrow (Q \rightarrow R)}{\rA}
  \proofstep{2}{S \rightarrow (R \rightarrow T)}{\rA}
  \proofstep{2}{P \rightarrow (S \rightarrow (R \rightarrow T))}{\FormulaRefAuto{a1i}{2}}
  \proofstep{1,2}{P \rightarrow (S \rightarrow (Q \rightarrow T))}{\FormulaRefAuto{syl5d}{1,3}}
\end{tabproof}

\FormulaThm[A nested syllogism inference with different antecedents]{%
  P \rightarrow (Q \rightarrow R) \dsep S \rightarrow (R \rightarrow T) \vdash S \rightarrow (P \rightarrow (Q \rightarrow T))%
}{syl9r}{%
  \DeltaRow{Propositions}{P\dsep Q\dsep R\dsep S\dsep T}%
}
\begin{tabproof}
  \proofstep{1}{P \rightarrow (Q \rightarrow R)}{\rA}
  \proofstep{2}{S \rightarrow (R \rightarrow T)}{\rA}
  \proofstep{1,2}{P \rightarrow (S \rightarrow (Q \rightarrow T))}{\FormulaRefAuto{syl9}{1,2}}
  \proofstep{1,2}{S \rightarrow (P \rightarrow (Q \rightarrow T))}{\FormulaRefAuto{com12}{3}}
\end{tabproof}

\FormulaThm[A nested syllogism inference]{%
  \begin{aligned}[t]
    &P \rightarrow (Q \rightarrow R)\dsep{}\\
    &P \rightarrow (Q \rightarrow (S \rightarrow T))\dsep{}\\
    &R \rightarrow (T \rightarrow U)\vdash{}\\
    &P \rightarrow (Q \rightarrow (S \rightarrow U))
  \end{aligned}
}{syl10}{%
  \DeltaRow{Propositions}{P\dsep Q\dsep R\dsep S\dsep T\dsep U}%
}
\begin{tabproof}
  \proofstep{1}{P \rightarrow (Q \rightarrow R)}{\rA}
  \proofstep{2}{P \rightarrow (Q \rightarrow (S \rightarrow T))}{\rA}
  \proofstep{3}{R \rightarrow (T \rightarrow U)}{\rA}
  \proofstep{1,3}{P \rightarrow (Q \rightarrow (T \rightarrow U))}{\FormulaRefAuto{syl6}{1,3}}
  \proofstep{1,2,3}{P \rightarrow (Q \rightarrow (S \rightarrow U))}{\FormulaRefAuto{syldd}{2,4}}
\end{tabproof}

\FormulaThm[Triple deduction introducing an antecedent to a wff]{%
  P \rightarrow (Q \rightarrow (R \rightarrow S)) \vdash P \rightarrow (Q \rightarrow (R \rightarrow (T \rightarrow S)))%
}{a1ddd}{%
  \DeltaRow{Propositions}{P\dsep Q\dsep R\dsep S\dsep T}%
}
\begin{tabproof}
  \proofstep{1}{P \rightarrow (Q \rightarrow (R \rightarrow S))}{\rA}
  \proofstep{}{S \rightarrow (T \rightarrow S)}{\FormulaRefAuto{ax-1}}
  \proofstep{1}{P \rightarrow (Q \rightarrow (R \rightarrow (T \rightarrow S)))}{\FormulaRefAuto{syl8}{1,2}}
\end{tabproof}

\FormulaThm[Deduction combining antecedents and consequents]{%
  P \rightarrow (Q \rightarrow R) \dsep P \rightarrow (S \rightarrow T) \vdash P \rightarrow ((R \rightarrow S) \rightarrow (Q \rightarrow T))%
}{imim12d}{%
  \DeltaRow{Propositions}{P\dsep Q\dsep R\dsep S\dsep T}%
}
\begin{tabproof}
  \proofstep{1}{P \rightarrow (Q \rightarrow R)}{\rA}
  \proofstep{2}{P \rightarrow (S \rightarrow T)}{\rA}
  \proofstep{2}{P \rightarrow ((R \rightarrow S) \rightarrow (R \rightarrow T))}{\FormulaRefAuto{imim2d}{2}}
  \proofstep{1,2}{P \rightarrow ((R \rightarrow S) \rightarrow (Q \rightarrow T))}{\FormulaRefAuto{syl5d}{1,3}}
\end{tabproof}

\FormulaThm[Deduction adding nested consequents]{%
  P \rightarrow (Q \rightarrow R) \vdash P \rightarrow ((R \rightarrow S) \rightarrow (Q \rightarrow S))%
}{imim1d}{%
  \DeltaRow{Propositions}{P\dsep Q\dsep R\dsep S}%
}
\begin{tabproof}
  \proofstep{1}{P \rightarrow (Q \rightarrow R)}{\rA}
  \proofstep{}{P \rightarrow (S \rightarrow S)}{\FormulaRefAuto{idd}}
  \proofstep{1}{P \rightarrow ((R \rightarrow S) \rightarrow (Q \rightarrow S))}{\FormulaRefAuto{imim12d}{1,2}}
\end{tabproof}

\FormulaThm[A closed form of syllogism]{%
  (P \rightarrow Q) \rightarrow ((Q \rightarrow R) \rightarrow (P \rightarrow R))%
}{imim1}{%
  \DeltaRow{Propositions}{P\dsep Q\dsep R}%
}
\begin{tabproof}
  \proofstep{}{(P \rightarrow Q) \rightarrow (P \rightarrow Q)}{\FormulaRefAuto{id}}
  \proofstep{}{(P \rightarrow Q) \rightarrow ((Q \rightarrow R) \rightarrow (P \rightarrow R))}{\FormulaRefAuto{imim1d}{1}}
\end{tabproof}

\FormulaThm[Theorem pm2.83]{%
  (P \rightarrow (Q \rightarrow R)) \rightarrow ((P \rightarrow (R \rightarrow S)) \rightarrow (P \rightarrow (Q \rightarrow S)))%
}{pm2.83}{%
  \DeltaRow{Propositions}{P\dsep Q\dsep R\dsep S}%
}
\begin{tabproof}
  \proofstep{}{(Q \rightarrow R) \rightarrow ((R \rightarrow S) \rightarrow (Q \rightarrow S))}{\FormulaRefAuto{imim1}}
  \proofstep{}{(P \rightarrow (Q \rightarrow R)) \rightarrow ((P \rightarrow (R \rightarrow S)) \rightarrow (P \rightarrow (Q \rightarrow S)))}{\FormulaRefAuto{imim3i}{1}}
\end{tabproof}

\FormulaThm[Peirce implies Roll]{%
  \begin{aligned}[t]
    &(((P \rightarrow Q) \rightarrow P) \rightarrow P)\rightarrow{}\\
    &\quad (((P \rightarrow Q) \rightarrow R)\rightarrow ((R \rightarrow P) \rightarrow P))
  \end{aligned}
}{peirceroll}{%
  \DeltaRow{Propositions}{P\dsep Q\dsep R}%
}
\begin{tabproof}
  \proofstep{}{%
    \begin{aligned}[t]
      &((P \rightarrow Q) \rightarrow R)\rightarrow{}\\
      &\quad ((R \rightarrow P) \rightarrow ((P \rightarrow Q) \rightarrow P))
    \end{aligned}
  }{\FormulaRefAuto{imim1}}

  \proofstep{}{%
    \begin{aligned}[t]
      &(((P \rightarrow Q) \rightarrow P) \rightarrow P)\rightarrow{}\\
      &\quad (((P \rightarrow Q) \rightarrow P) \rightarrow P)
    \end{aligned}
  }{\FormulaRefAuto{id}}

  \proofstep{}{%
    \begin{aligned}[t]
      &(((P \rightarrow Q) \rightarrow P) \rightarrow P)\rightarrow{}\\
      &\quad (((P \rightarrow Q) \rightarrow R)\rightarrow ((R \rightarrow P) \rightarrow P))
    \end{aligned}
  }{\FormulaRefAuto{syl9r}{1,2}}
\end{tabproof}

\FormulaThm[Commutation of antecedents. Swap 2nd and 3rd.]{%
  P \rightarrow (Q \rightarrow (R \rightarrow S)) \vdash P \rightarrow (R \rightarrow (Q \rightarrow S))%
}{com23}{%
  \DeltaRow{Propositions}{P\dsep Q\dsep R\dsep S}%
}
\begin{tabproof}
  \proofstep{1}{P \rightarrow (Q \rightarrow (R \rightarrow S))}{\rA}
  \proofstep{}{R \rightarrow ((R \rightarrow S) \rightarrow S)}{\FormulaRefAuto{pm2.27}}
  \proofstep{1}{P \rightarrow (R \rightarrow (Q \rightarrow S))}{\FormulaRefAuto{syl9}{1,2}}
\end{tabproof}

\FormulaThm[Commutation of antecedents. Rotate right.]{%
  P \rightarrow (Q \rightarrow (R \rightarrow S)) \vdash R \rightarrow (P \rightarrow (Q \rightarrow S))%
}{com3r}{%
  \DeltaRow{Propositions}{P\dsep Q\dsep R\dsep S}%
}
\begin{tabproof}
  \proofstep{1}{P \rightarrow (Q \rightarrow (R \rightarrow S))}{\rA}
  \proofstep{1}{P \rightarrow (R \rightarrow (Q \rightarrow S))}{\FormulaRefAuto{com23}{1}}
  \proofstep{1}{R \rightarrow (P \rightarrow (Q \rightarrow S))}{\FormulaRefAuto{com12}{2}}
\end{tabproof}

\FormulaThm[Commutation of antecedents. Swap 1st and 3rd.]{%
  P \rightarrow (Q \rightarrow (R \rightarrow S)) \vdash R \rightarrow (Q \rightarrow (P \rightarrow S))%
}{com13}{%
  \DeltaRow{Propositions}{P\dsep Q\dsep R\dsep S}%
}
\begin{tabproof}
  \proofstep{1}{P \rightarrow (Q \rightarrow (R \rightarrow S))}{\rA}
  \proofstep{1}{R \rightarrow (P \rightarrow (Q \rightarrow S))}{\FormulaRefAuto{com3r}{1}}
  \proofstep{1}{R \rightarrow (Q \rightarrow (P \rightarrow S))}{\FormulaRefAuto{com23}{2}}
\end{tabproof}

\FormulaThm[Commutation of antecedents. Rotate left.]{%
  P \rightarrow (Q \rightarrow (R \rightarrow S)) \vdash Q \rightarrow (R \rightarrow (P \rightarrow S))%
}{com3l}{%
  \DeltaRow{Propositions}{P\dsep Q\dsep R\dsep S}%
}
\begin{tabproof}
  \proofstep{1}{P \rightarrow (Q \rightarrow (R \rightarrow S))}{\rA}
  \proofstep{1}{R \rightarrow (P \rightarrow (Q \rightarrow S))}{\FormulaRefAuto{com3r}{1}}
  \proofstep{1}{Q \rightarrow (R \rightarrow (P \rightarrow S))}{\FormulaRefAuto{com3r}{2}}
\end{tabproof}

\FormulaThm[Swap antecedents]{%
  (P \rightarrow (Q \rightarrow R)) \rightarrow (Q \rightarrow (P \rightarrow R))%
}{pm2.04}{%
  \DeltaRow{Propositions}{P\dsep Q\dsep R}%
}
\begin{tabproof}
  \proofstep{}{(P \rightarrow (Q \rightarrow R)) \rightarrow (P \rightarrow (Q \rightarrow R))}{\FormulaRefAuto{id}}
  \proofstep{}{(P \rightarrow (Q \rightarrow R)) \rightarrow (Q \rightarrow (P \rightarrow R))}{\FormulaRefAuto{com23}{1}}
\end{tabproof}

\FormulaThm[Commutation of antecedents. Swap 3rd and 4th.]{%
  P \rightarrow (Q \rightarrow (R \rightarrow (S \rightarrow T))) \vdash P \rightarrow (Q \rightarrow (S \rightarrow (R \rightarrow T)))%
}{com34}{%
  \DeltaRow{Propositions}{P\dsep Q\dsep R\dsep S\dsep T}%
}
\begin{tabproof}
  \proofstep{1}{P \rightarrow (Q \rightarrow (R \rightarrow (S \rightarrow T)))}{\rA}
  \proofstep{}{(R \rightarrow (S \rightarrow T)) \rightarrow (S \rightarrow (R \rightarrow T))}{\FormulaRefAuto{pm2.04}}
  \proofstep{1}{P \rightarrow (Q \rightarrow (S \rightarrow (R \rightarrow T)))}{\FormulaRefAuto{syl6}{1,2}}
\end{tabproof}

\FormulaThm[Commutation of antecedents. Rotate left.]{%
  P \rightarrow (Q \rightarrow (R \rightarrow (S \rightarrow T))) \vdash Q \rightarrow (R \rightarrow (S \rightarrow (P \rightarrow T)))%
}{com4l}{%
  \DeltaRow{Propositions}{P\dsep Q\dsep R\dsep S\dsep T}%
}
\begin{tabproof}
  \proofstep{1}{P \rightarrow (Q \rightarrow (R \rightarrow (S \rightarrow T)))}{\rA}
  \proofstep{1}{Q \rightarrow (R \rightarrow (P \rightarrow (S \rightarrow T)))}{\FormulaRefAuto{com3l}{1}}
  \proofstep{1}{Q \rightarrow (R \rightarrow (S \rightarrow (P \rightarrow T)))}{\FormulaRefAuto{com34}{2}}
\end{tabproof}

\FormulaThm[Commutation of antecedents. Rotate twice.]{%
  P \rightarrow (Q \rightarrow (R \rightarrow (S \rightarrow T))) \vdash R \rightarrow (S \rightarrow (P \rightarrow (Q \rightarrow T)))%
}{com4t}{%
  \DeltaRow{Propositions}{P\dsep Q\dsep R\dsep S\dsep T}%
}
\begin{tabproof}
  \proofstep{1}{P \rightarrow (Q \rightarrow (R \rightarrow (S \rightarrow T)))}{\rA}
  \proofstep{1}{Q \rightarrow (R \rightarrow (S \rightarrow (P \rightarrow T)))}{\FormulaRefAuto{com4l}{1}}
  \proofstep{1}{R \rightarrow (S \rightarrow (P \rightarrow (Q \rightarrow T)))}{\FormulaRefAuto{com4l}{2}}
\end{tabproof}

\FormulaThm[Commutation of antecedents. Rotate right.]{%
  P \rightarrow (Q \rightarrow (R \rightarrow (S \rightarrow T))) \vdash S \rightarrow (P \rightarrow (Q \rightarrow (R \rightarrow T)))%
}{com4r}{%
  \DeltaRow{Propositions}{P\dsep Q\dsep R\dsep S\dsep T}%
}
\begin{tabproof}
  \proofstep{1}{P \rightarrow (Q \rightarrow (R \rightarrow (S \rightarrow T)))}{\rA}
  \proofstep{1}{R \rightarrow (S \rightarrow (P \rightarrow (Q \rightarrow T)))}{\FormulaRefAuto{com4t}{1}}
  \proofstep{1}{S \rightarrow (P \rightarrow (Q \rightarrow (R \rightarrow T)))}{\FormulaRefAuto{com4l}{2}}
\end{tabproof}

\FormulaThm[Commutation of antecedents. Swap 2nd and 4th.]{%
  P \rightarrow (Q \rightarrow (R \rightarrow (S \rightarrow T))) \vdash P \rightarrow (S \rightarrow (R \rightarrow (Q \rightarrow T)))%
}{com24}{%
  \DeltaRow{Propositions}{P\dsep Q\dsep R\dsep S\dsep T}%
}
\begin{tabproof}
  \proofstep{1}{P \rightarrow (Q \rightarrow (R \rightarrow (S \rightarrow T)))}{\rA}
  \proofstep{1}{R \rightarrow (S \rightarrow (P \rightarrow (Q \rightarrow T)))}{\FormulaRefAuto{com4t}{1}}
  \proofstep{1}{P \rightarrow (S \rightarrow (R \rightarrow (Q \rightarrow T)))}{\FormulaRefAuto{com13}{2}}
\end{tabproof}

\FormulaThm[Commutation of antecedents. Swap 1st and 4th.]{%
  P \rightarrow (Q \rightarrow (R \rightarrow (S \rightarrow T))) \vdash S \rightarrow (Q \rightarrow (R \rightarrow (P \rightarrow T)))%
}{com14}{%
  \DeltaRow{Propositions}{P\dsep Q\dsep R\dsep S\dsep T}%
}
\begin{tabproof}
  \proofstep{1}{P \rightarrow (Q \rightarrow (R \rightarrow (S \rightarrow T)))}{\rA}
  \proofstep{1}{Q \rightarrow (R \rightarrow (S \rightarrow (P \rightarrow T)))}{\FormulaRefAuto{com4l}{1}}
  \proofstep{1}{S \rightarrow (Q \rightarrow (R \rightarrow (P \rightarrow T)))}{\FormulaRefAuto{com3r}{2}}
\end{tabproof}

\FormulaThm[Commutation of antecedents. Swap 4th and 5th.]{%
  P \rightarrow (Q \rightarrow (R \rightarrow (S \rightarrow (T \rightarrow U)))) \vdash P \rightarrow (Q \rightarrow (R \rightarrow (T \rightarrow (S \rightarrow U))))%
}{com45}{%
  \DeltaRow{Propositions}{P\dsep Q\dsep R\dsep S\dsep T\dsep U}%
}
\begin{tabproof}
  \proofstep{1}{P \rightarrow (Q \rightarrow (R \rightarrow (S \rightarrow (T \rightarrow U))))}{\rA}
  \proofstep{}{(S \rightarrow (T \rightarrow U)) \rightarrow (T \rightarrow (S \rightarrow U))}{\FormulaRefAuto{pm2.04}}
  \proofstep{1}{P \rightarrow (Q \rightarrow (R \rightarrow (T \rightarrow (S \rightarrow U))))}{\FormulaRefAuto{syl8}{1,2}}
\end{tabproof}

\FormulaThm[Commutation of antecedents. Swap 3rd and 5th.]{%
  P \rightarrow (Q \rightarrow (R \rightarrow (S \rightarrow (T \rightarrow U)))) \vdash P \rightarrow (Q \rightarrow (T \rightarrow (S \rightarrow (R \rightarrow U))))%
}{com35}{%
  \DeltaRow{Propositions}{P\dsep Q\dsep R\dsep S\dsep T\dsep U}%
}
\begin{tabproof}
  \proofstep{1}{P \rightarrow (Q \rightarrow (R \rightarrow (S \rightarrow (T \rightarrow U))))}{\rA}
  \proofstep{1}{P \rightarrow (Q \rightarrow (S \rightarrow (R \rightarrow (T \rightarrow U))))}{\FormulaRefAuto{com34}{1}}
  \proofstep{1}{P \rightarrow (Q \rightarrow (S \rightarrow (T \rightarrow (R \rightarrow U))))}{\FormulaRefAuto{com45}{2}}
  \proofstep{1}{P \rightarrow (Q \rightarrow (T \rightarrow (S \rightarrow (R \rightarrow U))))}{\FormulaRefAuto{com34}{3}}
\end{tabproof}

\FormulaThm[Commutation of antecedents. Swap 2nd and 5th.]{%
  P \rightarrow (Q \rightarrow (R \rightarrow (S \rightarrow (T \rightarrow U)))) \vdash P \rightarrow (T \rightarrow (R \rightarrow (S \rightarrow (Q \rightarrow U))))%
}{com25}{%
  \DeltaRow{Propositions}{P\dsep Q\dsep R\dsep S\dsep T\dsep U}%
}
\begin{tabproof}
  \proofstep{1}{P \rightarrow (Q \rightarrow (R \rightarrow (S \rightarrow (T \rightarrow U))))}{\rA}
  \proofstep{1}{P \rightarrow (S \rightarrow (R \rightarrow (Q \rightarrow (T \rightarrow U))))}{\FormulaRefAuto{com24}{1}}
  \proofstep{1}{P \rightarrow (S \rightarrow (R \rightarrow (T \rightarrow (Q \rightarrow U))))}{\FormulaRefAuto{com45}{2}}
  \proofstep{1}{P \rightarrow (T \rightarrow (R \rightarrow (S \rightarrow (Q \rightarrow U))))}{\FormulaRefAuto{com24}{3}}
\end{tabproof}

\FormulaThm[Commutation of antecedents. Rotate left.]{%
  P \rightarrow (Q \rightarrow (R \rightarrow (S \rightarrow (T \rightarrow U)))) \vdash Q \rightarrow (R \rightarrow (S \rightarrow (T \rightarrow (P \rightarrow U))))%
}{com5l}{%
  \DeltaRow{Propositions}{P\dsep Q\dsep R\dsep S\dsep T\dsep U}%
}
\begin{tabproof}
  \proofstep{1}{P \rightarrow (Q \rightarrow (R \rightarrow (S \rightarrow (T \rightarrow U))))}{\rA}
  \proofstep{1}{Q \rightarrow (R \rightarrow (S \rightarrow (P \rightarrow (T \rightarrow U))))}{\FormulaRefAuto{com4l}{1}}
  \proofstep{1}{Q \rightarrow (R \rightarrow (S \rightarrow (T \rightarrow (P \rightarrow U))))}{\FormulaRefAuto{com45}{2}}
\end{tabproof}

\FormulaThm[Commutation of antecedents. Swap 1st and 5th.]{%
  \begin{aligned}[t]
    &P \rightarrow (Q \rightarrow (R \rightarrow (S \rightarrow (T \rightarrow U))))\vdash{}\\
    &\quad T \rightarrow (Q \rightarrow (R \rightarrow (S \rightarrow (P \rightarrow U))))
  \end{aligned}
}{com15}{%
  \DeltaRow{Propositions}{P\dsep Q\dsep R\dsep S\dsep T\dsep U}%
}
\begin{tabproof}
  \proofstep{1}{%
    \begin{aligned}[t]
      &P \rightarrow (Q \rightarrow (R \rightarrow (S \rightarrow (T \rightarrow U))))
    \end{aligned}
  }{\rA}
  \proofstep{1}{%
    \begin{aligned}[t]
      &Q \rightarrow (R \rightarrow (S \rightarrow (T \rightarrow (P \rightarrow U))))
    \end{aligned}
  }{\FormulaRefAuto{com5l}{1}}
  \proofstep{1}{%
    \begin{aligned}[t]
      &T \rightarrow (Q \rightarrow (R \rightarrow (S \rightarrow (P \rightarrow U))))
    \end{aligned}
  }{\FormulaRefAuto{com4r}{2}}
\end{tabproof}

\FormulaThm[Commutation of antecedents. Rotate left twice.]{%
  \begin{aligned}[t]
    &P \rightarrow (Q \rightarrow (R \rightarrow (S \rightarrow (T \rightarrow U))))\vdash{}\\
    &\quad R \rightarrow (S \rightarrow (T \rightarrow (P \rightarrow (Q \rightarrow U))))
  \end{aligned}
}{com52l}{%
  \DeltaRow{Propositions}{P\dsep Q\dsep R\dsep S\dsep T\dsep U}%
}
\begin{tabproof}
  \proofstep{1}{%
    \begin{aligned}[t]
      &P \rightarrow (Q \rightarrow (R \rightarrow (S \rightarrow (T \rightarrow U))))
    \end{aligned}
  }{\rA}
  \proofstep{1}{%
    \begin{aligned}[t]
      &Q \rightarrow (R \rightarrow (S \rightarrow (T \rightarrow (P \rightarrow U))))
    \end{aligned}
  }{\FormulaRefAuto{com5l}{1}}
  \proofstep{1}{%
    \begin{aligned}[t]
      &R \rightarrow (S \rightarrow (T \rightarrow (P \rightarrow (Q \rightarrow U))))
    \end{aligned}
  }{\FormulaRefAuto{com5l}{2}}
\end{tabproof}

\FormulaThm[Commutation of antecedents. Rotate right twice.]{%
  \begin{aligned}[t]
    &P \rightarrow (Q \rightarrow (R \rightarrow (S \rightarrow (T \rightarrow U))))\vdash{}\\
    &\quad S \rightarrow (T \rightarrow (P \rightarrow (Q \rightarrow (R \rightarrow U))))
  \end{aligned}
}{com52r}{%
  \DeltaRow{Propositions}{P\dsep Q\dsep R\dsep S\dsep T\dsep U}%
}
\begin{tabproof}
  \proofstep{1}{%
    \begin{aligned}[t]
      &P \rightarrow (Q \rightarrow (R \rightarrow (S \rightarrow (T \rightarrow U))))
    \end{aligned}
  }{\rA}
  \proofstep{1}{%
    \begin{aligned}[t]
      &R \rightarrow (S \rightarrow (T \rightarrow (P \rightarrow (Q \rightarrow U))))
    \end{aligned}
  }{\FormulaRefAuto{com52l}{1}}
  \proofstep{1}{%
    \begin{aligned}[t]
      &S \rightarrow (T \rightarrow (P \rightarrow (Q \rightarrow (R \rightarrow U))))
    \end{aligned}
  }{\FormulaRefAuto{com5l}{2}}
\end{tabproof}

\FormulaThm[Commutation of antecedents. Rotate right.]{%
  \begin{aligned}[t]
    &P \rightarrow (Q \rightarrow (R \rightarrow (S \rightarrow (T \rightarrow U))))\vdash{}\\
    &\quad T \rightarrow (P \rightarrow (Q \rightarrow (R \rightarrow (S \rightarrow U))))
  \end{aligned}
}{com5r}{%
  \DeltaRow{Propositions}{P\dsep Q\dsep R\dsep S\dsep T\dsep U}%
}
\begin{tabproof}
  \proofstep{1}{%
    \begin{aligned}[t]
      &P \rightarrow (Q \rightarrow (R \rightarrow (S \rightarrow (T \rightarrow U))))
    \end{aligned}
  }{\rA}
  \proofstep{1}{%
    \begin{aligned}[t]
      &R \rightarrow (S \rightarrow (T \rightarrow (P \rightarrow (Q \rightarrow U))))
    \end{aligned}
  }{\FormulaRefAuto{com52l}{1}}
  \proofstep{1}{%
    \begin{aligned}[t]
      &T \rightarrow (P \rightarrow (Q \rightarrow (R \rightarrow (S \rightarrow U))))
    \end{aligned}
  }{\FormulaRefAuto{com52l}{2}}
\end{tabproof}

\FormulaThm[Closed form of imim12i and of 3syl]{%
  ((P \rightarrow Q) \rightarrow ((R \rightarrow S) \rightarrow ((Q \rightarrow R) \rightarrow (P \rightarrow S))))
}{imim12}{%
  \DeltaRow{Propositions}{P\dsep Q\dsep R\dsep S}%
}
\begin{tabproof}
  \proofstep{}{%
    (R \rightarrow S) \rightarrow ((Q \rightarrow R) \rightarrow (Q \rightarrow S))
  }{\FormulaRefAuto{imim2}}
  \proofstep{}{%
    (P \rightarrow Q) \rightarrow ((Q \rightarrow S) \rightarrow (P \rightarrow S))
  }{\FormulaRefAuto{imim1}}
  \proofstep{}{%
    ((P \rightarrow Q) \rightarrow ((R \rightarrow S) \rightarrow ((Q \rightarrow R) \rightarrow (P \rightarrow S))))
  }{\FormulaRefAuto{syl9r}{1,2}}
\end{tabproof}

\FormulaThm[Elimination of a nested antecedent]{%
  (((P \rightarrow Q) \rightarrow R) \rightarrow (Q \rightarrow R))
}{jarr}{%
  \DeltaRow{Propositions}{P\dsep Q\dsep R}%
}
\begin{tabproof}
  \proofstep{}{Q \rightarrow (P \rightarrow Q)}{\FormulaRefAuto{ax-1}}
  \proofstep{}{(((P \rightarrow Q) \rightarrow R) \rightarrow (Q \rightarrow R))}{\FormulaRefAuto{imim1i}{1}}
\end{tabproof}

\begin{remark}
\FormulaRefAuto{jarr} ist eine Kürzungsregel: der zusätzliche Antezedens \((P \rightarrow Q)\) verschwindet.
\end{remark}

\FormulaThm[Inference associated with jarr]{%
  ((P \rightarrow Q) \rightarrow R) \vdash Q \rightarrow R
}{jarri}{%
  \DeltaRow{Propositions}{P\dsep Q\dsep R}%
}
\begin{tabproof}
  \proofstep{}{Q \rightarrow (P \rightarrow Q)}{\FormulaRefAuto{ax-1}}
  \proofstep{2}{(P \rightarrow Q) \rightarrow R}{\rA}
  \proofstep{2}{Q \rightarrow R}{\FormulaRefAuto{syl}{1,2}}
\end{tabproof}

\FormulaThm[Deduction associated with pm2.86]{%
  P \rightarrow ((Q \rightarrow R) \rightarrow (Q \rightarrow S)) \vdash P \rightarrow (Q \rightarrow (R \rightarrow S))
}{pm2.86d}{%
  \DeltaRow{Propositions}{P\dsep Q\dsep R\dsep S}%
}
\begin{tabproof}
  \proofstep{}{R \rightarrow (Q \rightarrow R)}{\FormulaRefAuto{ax-1}}
  \proofstep{2}{P \rightarrow ((Q \rightarrow R) \rightarrow (Q \rightarrow S))}{\rA}
  \proofstep{2}{P \rightarrow (R \rightarrow (Q \rightarrow S))}{\FormulaRefAuto{syl5}{1,2}}
  \proofstep{2}{P \rightarrow (Q \rightarrow (R \rightarrow S))}{\FormulaRefAuto{com23}{3}}
\end{tabproof}

\FormulaThm[Converse of Axiom ax-2]{%
  (((P \rightarrow Q) \rightarrow (P \rightarrow R)) \rightarrow (P \rightarrow (Q \rightarrow R)))
}{pm2.86}{%
  \DeltaRow{Propositions}{P\dsep Q\dsep R}%
}
\begin{tabproof}
  \proofstep{}{%
    ((P \rightarrow Q) \rightarrow (P \rightarrow R))
    \rightarrow
    ((P \rightarrow Q) \rightarrow (P \rightarrow R))
  }{\FormulaRefAuto{id}}
  \proofstep{}{%
    (((P \rightarrow Q) \rightarrow (P \rightarrow R)) \rightarrow (P \rightarrow (Q \rightarrow R)))
  }{\FormulaRefAuto{pm2.86d}{1}}
\end{tabproof}

\begin{remark}
\FormulaRefAuto{pm2.86} ist die Umkehrung von \FormulaRefAuto{ax-2}: ein gemeinsames Antezedens wird wieder in eine verschachtelte Implikation eingezogen.
\end{remark}

\FormulaThm[Inference associated with pm2.86]{%
  ((P \rightarrow Q) \rightarrow (P \rightarrow R)) \vdash P \rightarrow (Q \rightarrow R)
}{pm2.86i}{%
  \DeltaRow{Propositions}{P\dsep Q\dsep R}%
}
\begin{tabproof}
  \proofstep{1}{(P \rightarrow Q) \rightarrow (P \rightarrow R)}{\rA}
  \proofstep{1}{Q \rightarrow (P \rightarrow R)}{\FormulaRefAuto{jarri}{1}}
  \proofstep{1}{P \rightarrow (Q \rightarrow R)}{\FormulaRefAuto{com12}{2}}
\end{tabproof}

\FormulaThm[The Linearity Axiom of Lukasiewicz]{%
  (((P \rightarrow Q) \rightarrow (Q \rightarrow P)) \rightarrow (Q \rightarrow P))
}{loolin}{%
  \DeltaRow{Propositions}{P\dsep Q}%
}
\begin{tabproof}
  \proofstep{}{%
    (((P \rightarrow Q) \rightarrow (Q \rightarrow P)) \rightarrow (Q \rightarrow (Q \rightarrow P)))
  }{\FormulaRefAuto{jarr}}
  \proofstep{}{%
    (((P \rightarrow Q) \rightarrow (Q \rightarrow P)) \rightarrow (Q \rightarrow P))
  }{\FormulaRefAuto{pm2.43d}{1}}
\end{tabproof}
\begin{remark}
\FormulaRefAuto{loolin} gehört nicht zum eigentlichen Kern der hier benutzten Axiomatik, sondern erscheint als in unserem Kalkül beweisbare Form eines Axioms der unendlichwertigen Aussagenlogik von Łukasiewicz. Der Satz ist deshalb besonders als Vergleich zwischen verschiedenen Aussagenkalkülen aufschlussreich.
\end{remark}

\FormulaThm[An alternate for the Linearity Axiom of Lukasiewicz]{%
  (((P \rightarrow Q) \rightarrow (P \rightarrow R)) \rightarrow ((Q \rightarrow P) \rightarrow (Q \rightarrow R)))
}{loowoz}{%
  \DeltaRow{Propositions}{P\dsep Q\dsep R}%
}
\begin{tabproof}
  \proofstep{}{%
    (((P \rightarrow Q) \rightarrow (P \rightarrow R)) \rightarrow (Q \rightarrow (P \rightarrow R)))
  }{\FormulaRefAuto{jarr}}
  \proofstep{}{%
    (((P \rightarrow Q) \rightarrow (P \rightarrow R)) \rightarrow ((Q \rightarrow P) \rightarrow (Q \rightarrow R)))
  }{\FormulaRefAuto{a2d}{1}}
\end{tabproof}

\section{Logical negation}


\FormulaThm[Alias for ax-3]{%
  ((\neg P \rightarrow \neg Q) \rightarrow (Q \rightarrow P))
}{con4}{%
  \DeltaRow{Propositions}{P\dsep Q}%
}
\begin{tabproof}
  \proofstep{}{(\neg P \rightarrow \neg Q) \rightarrow (Q \rightarrow P)}{\FormulaRefAuto{ax-3}}
\end{tabproof}
\begin{remark}
\FormulaRefAuto{con4} ist die Negationsform von \FormulaRefAuto{ax-3} und markiert den eigentlichen Einstieg in die klassische Negationslogik.
\end{remark}

\FormulaThm[Inference associated with con4]{%
  (\neg P \rightarrow \neg Q) \vdash Q \rightarrow P
}{con4i}{%
  \DeltaRow{Propositions}{P\dsep Q}%
}
\begin{tabproof}
  \proofstep{1}{\neg P \rightarrow \neg Q}{\rA}
  \proofstep{}{(\neg P \rightarrow \neg Q) \rightarrow (Q \rightarrow P)}{\FormulaRefAuto{con4}}
  \proofstep{1}{Q \rightarrow P}{\FormulaRefAuto{ax-mp}{1,2}}
\end{tabproof}
\begin{remark}
Metamath vermerkt zu \FormulaRefAuto{con4i} einen kürzeren Beweis über \FormulaRefAuto{notnot} und \FormulaRefAuto{nsyl2}. Der hier naheliegende Zugang über \FormulaRefAuto{con4} ist jedoch für den Aufbau des Abschnitts günstiger, weil er ohne spätere Sätze über die Doppelnegation auskommt.
\end{remark}


\FormulaThm[Deduction associated with con4]{%
  P \rightarrow (\neg Q \rightarrow \neg R) \vdash P \rightarrow (R \rightarrow Q)
}{con4d}{%
  \DeltaRow{Propositions}{P\dsep Q\dsep R}%
}
\begin{tabproof}
  \proofstep{1}{P \rightarrow (\neg Q \rightarrow \neg R)}{\rA}
  \proofstep{}{(\neg Q \rightarrow \neg R) \rightarrow (R \rightarrow Q)}{\FormulaRefAuto{con4}}
  \proofstep{1}{P \rightarrow (R \rightarrow Q)}{\FormulaRefAuto{syl}{1,2}}
\end{tabproof}

\FormulaThm[The rule of modus tollens]{%
  P \dsep \neg Q \rightarrow \neg P \vdash Q
}{mt4}{%
  \DeltaRow{Propositions}{P\dsep Q}%
}
\begin{tabproof}
  \proofstep{1}{P}{\rA}
  \proofstep{2}{\neg Q \rightarrow \neg P}{\rA}
  \proofstep{2}{P \rightarrow Q}{\FormulaRefAuto{con4i}{2}}
  \proofstep{1,2}{Q}{\FormulaRefAuto{ax-mp}{1,3}}
\end{tabproof}

\FormulaThm[Modus tollens deduction]{%
  P \rightarrow Q \dsep P \rightarrow (\neg R \rightarrow \neg Q) \vdash P \rightarrow R
}{mt4d}{%
  \DeltaRow{Propositions}{P\dsep Q\dsep R}%
}
\begin{tabproof}
  \proofstep{1}{P \rightarrow Q}{\rA}
  \proofstep{2}{P \rightarrow (\neg R \rightarrow \neg Q)}{\rA}
  \proofstep{2}{P \rightarrow (Q \rightarrow R)}{\FormulaRefAuto{con4d}{2}}
  \proofstep{1,2}{P \rightarrow R}{\FormulaRefAuto{mpd}{1,3}}
\end{tabproof}

\FormulaThm[Modus tollens inference]{%
  R \dsep P \rightarrow (\neg Q \rightarrow \neg R) \vdash P \rightarrow Q
}{mt4i}{%
  \DeltaRow{Propositions}{P\dsep Q\dsep R}%
}
\begin{tabproof}
  \proofstep{1}{R}{\rA}
  \proofstep{1}{P \rightarrow R}{\FormulaRefAuto{a1i}{1}}
  \proofstep{2}{P \rightarrow (\neg Q \rightarrow \neg R)}{\rA}
  \proofstep{1,2}{P \rightarrow Q}{\FormulaRefAuto{mt4d}{1,3}}
\end{tabproof}

\FormulaThm[A contradiction implies anything]{%
  \neg P \vdash P \rightarrow Q
}{pm2.21i}{%
  \DeltaRow{Propositions}{P\dsep Q}%
}
\begin{tabproof}
  \proofstep{1}{\neg P}{\rA}
  \proofstep{1}{\neg Q \rightarrow \neg P}{\FormulaRefAuto{a1i}{1}}
  \proofstep{1}{P \rightarrow Q}{\FormulaRefAuto{con4i}{2}}
\end{tabproof}
\begin{remark}
\FormulaRefAuto{pm2.21} ist das Duns-Scotus-Gesetz: aus einem Widerspruch folgt jede Aussage.
\end{remark}

\FormulaThm[A contradiction implies anything]{%
  P \dsep \neg P \vdash Q
}{pm2.24ii}{%
  \DeltaRow{Propositions}{P\dsep Q}%
}
\begin{tabproof}
  \proofstep{1}{P}{\rA}
  \proofstep{2}{\neg P}{\rA}
  \proofstep{2}{P \rightarrow Q}{\FormulaRefAuto{pm2.21i}{2}}
  \proofstep{1,2}{Q}{\FormulaRefAuto{ax-mp}{1,3}}
\end{tabproof}

\FormulaThm[A contradiction implies anything]{%
  P \rightarrow \neg Q \vdash P \rightarrow (Q \rightarrow R)
}{pm2.21d}{%
  \DeltaRow{Propositions}{P\dsep Q\dsep R}%
}
\begin{tabproof}
  \proofstep{1}{P \rightarrow \neg Q}{\rA}
  \proofstep{1}{P \rightarrow (\neg R \rightarrow \neg Q)}{\FormulaRefAuto{a1d}{1}}
  \proofstep{1}{P \rightarrow (Q \rightarrow R)}{\FormulaRefAuto{con4d}{2}}
\end{tabproof}

\FormulaThm[Alternate proof of pm2.21dd]{%
  P \rightarrow Q \dsep P \rightarrow \neg Q \vdash P \rightarrow R
}{pm2.21ddALT}{%
  \DeltaRow{Propositions}{P\dsep Q\dsep R}%
}
\begin{tabproof}
  \proofstep{1}{P \rightarrow Q}{\rA}
  \proofstep{2}{P \rightarrow \neg Q}{\rA}
  \proofstep{2}{P \rightarrow (Q \rightarrow R)}{\FormulaRefAuto{pm2.21d}{2}}
  \proofstep{1,2}{P \rightarrow R}{\FormulaRefAuto{mpd}{1,3}}
\end{tabproof}

\FormulaThm[From a wff and its negation, anything follows]{%
  (\neg P \rightarrow (P \rightarrow Q))
}{pm2.21}{%
  \DeltaRow{Propositions}{P\dsep Q}%
}
\begin{tabproof}
  \proofstep{}{\neg P \rightarrow \neg P}{\FormulaRefAuto{id}}
  \proofstep{}{(\neg P \rightarrow (P \rightarrow Q))}{\FormulaRefAuto{pm2.21d}{1}}
\end{tabproof}

\FormulaThm[Commuted form of pm2.21]{%
  (P \rightarrow (\neg P \rightarrow Q))
}{pm2.24}{%
  \DeltaRow{Propositions}{P\dsep Q}%
}
\begin{tabproof}
  \proofstep{}{(\neg P \rightarrow (P \rightarrow Q))}{\FormulaRefAuto{pm2.21}}
  \proofstep{}{(P \rightarrow (\neg P \rightarrow Q))}{\FormulaRefAuto{com12}{1}}
\end{tabproof}

\FormulaThm[Elimination of a nested antecedent]{%
  (((P \rightarrow Q) \rightarrow R) \rightarrow (\neg P \rightarrow R))
}{jarl}{%
  \DeltaRow{Propositions}{P\dsep Q\dsep R}%
}
\begin{tabproof}
  \proofstep{}{(\neg P \rightarrow (P \rightarrow Q))}{\FormulaRefAuto{pm2.21}}
  \proofstep{}{(((P \rightarrow Q) \rightarrow R) \rightarrow (\neg P \rightarrow R))}{\FormulaRefAuto{imim1i}{1}}
\end{tabproof}
\begin{remark}
\FormulaRefAuto{jarl} ist das negationslogische Gegenstück zu \FormulaRefAuto{jarr}. Statt aus \(Q\) wird hier aus \(\neg P\) auf \(R\) geschlossen. Der Satz ist nützlich, wenn ein verschachtelter Vordersatz \((P \rightarrow Q)\) unter einer Widerlegungsannahme über \(P\) beseitigt werden soll.
\end{remark}

\FormulaThm[Inference associated with jarl]{%
  ((P \rightarrow Q) \rightarrow R) \vdash \neg P \rightarrow R
}{jarli}{%
  \DeltaRow{Propositions}{P\dsep Q\dsep R}%
}
\begin{tabproof}
  \proofstep{}{(\neg P \rightarrow (P \rightarrow Q))}{\FormulaRefAuto{pm2.21}}
  \proofstep{1}{(P \rightarrow Q) \rightarrow R}{\rA}
  \proofstep{1}{\neg P \rightarrow R}{\FormulaRefAuto{syl}{1,2}}
\end{tabproof}

\FormulaThm[Deduction form of the Clavius law]{%
  P \rightarrow (\neg Q \rightarrow Q) \vdash P \rightarrow Q
}{pm2.18d}{%
  \DeltaRow{Propositions}{P\dsep Q}%
}
\begin{tabproof}
  \proofstep{}{P \rightarrow P}{\FormulaRefAuto{id}}
  \proofstep{2}{P \rightarrow (\neg Q \rightarrow Q)}{\rA}
  \proofstep{}{(\neg Q \rightarrow (Q \rightarrow \neg P))}{\FormulaRefAuto{pm2.21}}
  \proofstep{2}{P \rightarrow (\neg Q \rightarrow \neg P)}{\FormulaRefAuto{sylcom}{2,3}}
  \proofstep{2}{P \rightarrow Q}{\FormulaRefAuto{mt4d}{1,4}}
\end{tabproof}

\FormulaThm[Clavius law]{%
  ((\neg P \rightarrow P) \rightarrow P)
}{pm2.18}{%
  \DeltaRow{Propositions}{P}%
}
\begin{tabproof}
  \proofstep{}{(\neg P \rightarrow P) \rightarrow (\neg P \rightarrow P)}{\FormulaRefAuto{id}}
  \proofstep{}{((\neg P \rightarrow P) \rightarrow P)}{\FormulaRefAuto{pm2.18d}{1}}
\end{tabproof}
\begin{remark}
\FormulaRefAuto{pm2.18} ist das Clavius-Gesetz und rechtfertigt Beweise durch Widerspruch.
\end{remark}

\FormulaThm[Inference associated with the Clavius law]{%
  (\neg P \rightarrow P) \vdash P
}{pm2.18i}{%
  \DeltaRow{Propositions}{P}%
}
\begin{tabproof}
  \proofstep{1}{\neg P \rightarrow P}{\rA}
  \proofstep{}{(\neg P \rightarrow P) \rightarrow P}{\FormulaRefAuto{pm2.18}}
  \proofstep{1}{P}{\FormulaRefAuto{ax-mp}{1,2}}
\end{tabproof}

\FormulaThm[Double negation elimination]{%
  (\neg \neg P \rightarrow P)
}{notnotr}{%
  \DeltaRow{Propositions}{P}%
}
\begin{tabproof}
  \proofstep{}{((\neg P \rightarrow P) \rightarrow P)}{\FormulaRefAuto{pm2.18}}
  \proofstep{}{(\neg \neg P \rightarrow P)}{\FormulaRefAuto{jarli}{1}}
\end{tabproof}

\FormulaThm[Inference associated with notnotr]{%
  \neg \neg P \vdash P
}{notnotri}{%
  \DeltaRow{Propositions}{P}%
}
\begin{tabproof}
  \proofstep{1}{\neg \neg P}{\rA}
  \proofstep{1}{\neg P \rightarrow \neg \neg \neg P}{\FormulaRefAuto{pm2.21i}{1}}
  \proofstep{1}{P}{\FormulaRefAuto{mt4}{1,2}}
\end{tabproof}
\begin{remark}
\FormulaRefAuto{notnotr} ist klassisch, aber nicht intuitionistisch allgemein gültig.
\end{remark}

\FormulaThm[Alternate proof of notnotri]{%
  \neg \neg P \vdash P
}{notnotriALT}{%
  \DeltaRow{Propositions}{P}%
}
\begin{tabproof}
  \proofstep{1}{\neg \neg P}{\rA}
  \proofstep{1}{\neg P \rightarrow P}{\FormulaRefAuto{pm2.21i}{1}}
  \proofstep{1}{P}{\FormulaRefAuto{pm2.18i}{2}}
\end{tabproof}

\FormulaThm[Deduction associated with notnotr]{%
  P \rightarrow \neg \neg Q \vdash P \rightarrow Q
}{notnotrd}{%
  \DeltaRow{Propositions}{P\dsep Q}%
}
\begin{tabproof}
  \proofstep{1}{P \rightarrow \neg \neg Q}{\rA}
  \proofstep{}{(\neg \neg Q \rightarrow Q)}{\FormulaRefAuto{notnotr}}
  \proofstep{1}{P \rightarrow Q}{\FormulaRefAuto{syl}{1,2}}
\end{tabproof}

\FormulaThm[A contraposition deduction]{%
  P \rightarrow (Q \rightarrow \neg R) \vdash P \rightarrow (R \rightarrow \neg Q)
}{con2d}{%
  \DeltaRow{Propositions}{P\dsep Q\dsep R}%
}
\begin{tabproof}
  \proofstep{}{(\neg \neg Q \rightarrow Q)}{\FormulaRefAuto{notnotr}}
  \proofstep{2}{P \rightarrow (Q \rightarrow \neg R)}{\rA}
  \proofstep{2}{P \rightarrow (\neg \neg Q \rightarrow \neg R)}{\FormulaRefAuto{syl5}{1,2}}
  \proofstep{2}{P \rightarrow (R \rightarrow \neg Q)}{\FormulaRefAuto{con4d}{3}}
\end{tabproof}

\FormulaThm[Contraposition]{%
  ((P \rightarrow \neg Q) \rightarrow (Q \rightarrow \neg P))
}{con2}{%
  \DeltaRow{Propositions}{P\dsep Q}%
}
\begin{tabproof}
  \proofstep{}{(P \rightarrow \neg Q) \rightarrow (P \rightarrow \neg Q)}{\FormulaRefAuto{id}}
  \proofstep{}{((P \rightarrow \neg Q) \rightarrow (Q \rightarrow \neg P))}{\FormulaRefAuto{con2d}{1}}
\end{tabproof}
\begin{remark}
\FormulaRefAuto{con2} ist die klassische Kontraposition in reiner Implikationsform. Viele spätere Negationsschlüsse, insbesondere \FormulaRefAuto{mt2d}, \FormulaRefAuto{nsyl} und \FormulaRefAuto{nsyl2}, lassen sich als Anwendungen genau dieses Grundmusters lesen.
\end{remark}

\FormulaThm[Modus tollens deduction]{%
  P \rightarrow R \dsep P \rightarrow (Q \rightarrow \neg R) \vdash P \rightarrow \neg Q
}{mt2d}{%
  \DeltaRow{Propositions}{P\dsep Q\dsep R}%
}
\begin{tabproof}
  \proofstep{1}{P \rightarrow R}{\rA}
  \proofstep{2}{P \rightarrow (Q \rightarrow \neg R)}{\rA}
  \proofstep{2}{P \rightarrow (R \rightarrow \neg Q)}{\FormulaRefAuto{con2d}{2}}
  \proofstep{1,2}{P \rightarrow \neg Q}{\FormulaRefAuto{mpd}{1,3}}
\end{tabproof}

\FormulaThm[Modus tollens inference]{%
  R \dsep P \rightarrow (Q \rightarrow \neg R) \vdash P \rightarrow \neg Q
}{mt2i}{%
  \DeltaRow{Propositions}{P\dsep Q\dsep R}%
}
\begin{tabproof}
  \proofstep{1}{R}{\rA}
  \proofstep{1}{P \rightarrow R}{\FormulaRefAuto{a1i}{1}}
  \proofstep{2}{P \rightarrow (Q \rightarrow \neg R)}{\rA}
  \proofstep{1,2}{P \rightarrow \neg Q}{\FormulaRefAuto{mt2d}{2,3}}
\end{tabproof}

\FormulaThm[A negated syllogism inference]{%
  P \rightarrow \neg Q \dsep R \rightarrow Q \vdash R \rightarrow \neg P
}{nsyl3}{%
  \DeltaRow{Propositions}{P\dsep Q\dsep R}%
}
\begin{tabproof}
  \proofstep{1}{R \rightarrow Q}{\rA}
  \proofstep{2}{P \rightarrow \neg Q}{\rA}
  \proofstep{2}{R \rightarrow (P \rightarrow \neg Q)}{\FormulaRefAuto{a1i}{2}}
  \proofstep{1,2}{R \rightarrow \neg P}{\FormulaRefAuto{mt2d}{1,3}}
\end{tabproof}

\FormulaThm[A contraposition inference]{%
  P \rightarrow \neg Q \vdash Q \rightarrow \neg P
}{con2i}{%
  \DeltaRow{Propositions}{P\dsep Q}%
}
\begin{tabproof}
  \proofstep{1}{P \rightarrow \neg Q}{\rA}
  \proofstep{}{Q \rightarrow Q}{\FormulaRefAuto{id}}
  \proofstep{1}{Q \rightarrow \neg P}{\FormulaRefAuto{nsyl3}{1,2}}
\end{tabproof}

\FormulaThm[A negated syllogism inference]{%
  P \rightarrow \neg Q \dsep R \rightarrow Q \vdash P \rightarrow \neg R
}{nsyl}{%
  \DeltaRow{Propositions}{P\dsep Q\dsep R}%
}
\begin{tabproof}
  \proofstep{1}{P \rightarrow \neg Q}{\rA}
  \proofstep{2}{R \rightarrow Q}{\rA}
  \proofstep{1,2}{R \rightarrow \neg P}{\FormulaRefAuto{nsyl3}{1,2}}
  \proofstep{1,2}{P \rightarrow \neg R}{\FormulaRefAuto{con2i}{3}}
\end{tabproof}

\FormulaThm[A negated syllogism inference]{%
  P \rightarrow \neg Q \dsep \neg R \rightarrow Q \vdash P \rightarrow R
}{nsyl2}{%
  \DeltaRow{Propositions}{P\dsep Q\dsep R}%
}
\begin{tabproof}
  \proofstep{1}{P \rightarrow \neg Q}{\rA}
  \proofstep{2}{\neg R \rightarrow Q}{\rA}
  \proofstep{1,2}{\neg R \rightarrow \neg P}{\FormulaRefAuto{nsyl3}{1,2}}
  \proofstep{1,2}{P \rightarrow R}{\FormulaRefAuto{con4i}{3}}
\end{tabproof}

\FormulaThm[Double negation introduction]{%
  (P \rightarrow \neg \neg P)
}{notnot}{%
  \DeltaRow{Propositions}{P}%
}
\begin{tabproof}
  \proofstep{}{\neg P \rightarrow \neg P}{\FormulaRefAuto{id}}
  \proofstep{}{P \rightarrow \neg \neg P}{\FormulaRefAuto{con2i}{1}}
\end{tabproof}
\begin{remark}
Zusammen mit \FormulaRefAuto{notnotr} liefert \FormulaRefAuto{notnot} die in der klassischen Logik gültige Äquivalenz zwischen \(P\) und \(\neg\neg P\). Im Unterschied zu \FormulaRefAuto{notnotr} ist diese Einführungsrichtung auch intuitionistisch unproblematisch.
\end{remark}

\FormulaThm[Inference associated with notnot]{%
  P \vdash \neg \neg P
}{notnoti}{%
  \DeltaRow{Propositions}{P}%
}
\begin{tabproof}
  \proofstep{1}{P}{\rA}
  \proofstep{}{P \rightarrow \neg \neg P}{\FormulaRefAuto{notnot}}
  \proofstep{1}{\neg \neg P}{\FormulaRefAuto{ax-mp}{1,2}}
\end{tabproof}

\FormulaThm[Deduction associated with notnot]{%
  P \rightarrow Q \vdash P \rightarrow \neg \neg Q
}{notnotd}{%
  \DeltaRow{Propositions}{P\dsep Q}%
}
\begin{tabproof}
  \proofstep{1}{P \rightarrow Q}{\rA}
  \proofstep{}{Q \rightarrow \neg \neg Q}{\FormulaRefAuto{notnot}}
  \proofstep{1}{P \rightarrow \neg \neg Q}{\FormulaRefAuto{syl}{1,2}}
\end{tabproof}


\PrintFormulaStackById

\end{document}